% Options for packages loaded elsewhere
\PassOptionsToPackage{unicode}{hyperref}
\PassOptionsToPackage{hyphens}{url}
\documentclass[
]{article}
\usepackage{xcolor}
\usepackage{amsmath,amssymb}
\setcounter{secnumdepth}{-\maxdimen} % remove section numbering
\usepackage{iftex}
\ifPDFTeX
  \usepackage[T1]{fontenc}
  \usepackage[utf8]{inputenc}
  \usepackage{textcomp} % provide euro and other symbols
\else % if luatex or xetex
  \usepackage{unicode-math} % this also loads fontspec
  \defaultfontfeatures{Scale=MatchLowercase}
  \defaultfontfeatures[\rmfamily]{Ligatures=TeX,Scale=1}
\fi
\usepackage{lmodern}
\ifPDFTeX\else
  % xetex/luatex font selection
\fi
% Use upquote if available, for straight quotes in verbatim environments
\IfFileExists{upquote.sty}{\usepackage{upquote}}{}
\IfFileExists{microtype.sty}{% use microtype if available
  \usepackage[]{microtype}
  \UseMicrotypeSet[protrusion]{basicmath} % disable protrusion for tt fonts
}{}
\makeatletter
\@ifundefined{KOMAClassName}{% if non-KOMA class
  \IfFileExists{parskip.sty}{%
    \usepackage{parskip}
  }{% else
    \setlength{\parindent}{0pt}
    \setlength{\parskip}{6pt plus 2pt minus 1pt}}
}{% if KOMA class
  \KOMAoptions{parskip=half}}
\makeatother
\setlength{\emergencystretch}{3em} % prevent overfull lines
\providecommand{\tightlist}{%
  \setlength{\itemsep}{0pt}\setlength{\parskip}{0pt}}
\usepackage{bookmark}
\IfFileExists{xurl.sty}{\usepackage{xurl}}{} % add URL line breaks if available
\urlstyle{same}
\hypersetup{
  hidelinks,
  pdfcreator={LaTeX via pandoc}}

\author{}
\date{}

\begin{document}

\textbf{Ratushnyi Neural Repair Coding}

\emph{A Framework for Lossless Neural Compression via Side‑Information
Repair: Stream, Positional, and Tensor Entropy‑Field Coding}

Pavlo Ratushnyi

\emph{Independent Research --- Draft v0.8, May 2026}

\textbf{Abstract}

We introduce Ratushnyi Neural Repair Coding (RNR Coding), a framework
for lossless compression in which a deterministic,
decoder-reproducible neural or hybrid model produces side information
about a source block, and a short repair description reconstructs the
exact original from that side information. The framework is organized
along the dimensionality and target of prediction. Type-I performs
one-dimensional stream repair; Type-II performs two-dimensional
positional symbol coding over a partition of positions among symbol
classes; Type-III performs multi-dimensional tensor entropy-field
coding, divided into III-A (entropy-field meta-prediction), III-B
(byte-tensor coding of joint nibble structure), and III-C
(token-tensor coding over a learned hierarchically indexed
dictionary). For each subtype we give a concrete construction and an
achievability bound; a matching converse is given only for Type-I
(Theorem 4.2). Complexity results include NP-hardness of the
varying-parameter learned-code labeling problem (Theorem 4.3) and of
a restricted normal-form dictionary-only case of Type-III-C (Theorem
6.7), plus an APX-hardness transfer for the same restricted case
with an inherited constant $\gamma_{\mathrm{SGP}}>1$ (Theorem 6.8);
hardness of the unrestricted joint (model, dictionary, repair)
problem and the analogous inapproximability for $E_{12}$ are open. Worked numerical examples are
given for Type-I, Type-II, and Type-III-B; III-A and III-C are
illustrated through their composition and dictionary-construction
sections respectively. Two structural results identify the
information-theoretic quantity governing factorized-vs-joint
separation: a $\Theta(N)$ amortized separation for a specific
linear-code family in semi-non-factorized Type-II (Theorem 5.4), and
the identity that for an arbitrary source $D_N$ the gap equals
multi-information $\mathrm{TC}(D_N)$ (Theorem 5.5, where the gap may
vanish on sources without joint structure). Type-III-A admits a tight
$2\varepsilon$ approximate-oracle bound (Theorem 6.1, with matching
lower bound Theorem 6.2). The framework's deployment-relevant feature is random-access decoding
with cost
$O(\log(|X|/K) \cdot \log L(R) + (K + k) \cdot \mathrm{cost}(M))$ for
$W$-bounded predictors (Theorem 10.3), where $|X|$ is the
uncompressed source length --- polylogarithmic in archive size
$|X|$ once the sync-point spacing $K$, read window $k$, predictor
cost, and $\log L(R)$ are treated as fixed (or polylogarithmic in
$|X|$) parameters. Under an additional probability-floor assumption
on the predictor (Section 10.7), the sync-point compression overhead
per sync point is bounded by $W \cdot \log_2(1/\eta)$ where $\eta$
is the floor (a sharper bound subtracts a warm-context rate term
under an additional uniform-warm-rate assumption); we conjecture the
actual overhead is on the order of one percent for natural data at
typical parameters, but neither a tight upper bound in general nor
a measured value on a specific corpus is provided here. The
ratio comparison against existing dictionary-based seekable formats
is an empirical conjecture (Section 11), not a theorem. Bit-exact cross-platform reproduction is
reduced to a precision-and-determinism specification (Theorem 10.1);
published integer-arithmetic methods (integer SGD, signSGD, integer
quantized inference) provide candidate components that an implementer
must verify against the full specification on a chosen target. We do not claim universal
dominance --- no lossless code can map every input of length $n$ to a
strictly shorter one --- and the framework is evaluated conditionally
on data classes for which side information reduces residual
uncertainty by more than the model and repair overhead costs.

\textbf{Keywords:} lossless compression, neural compression, side
information, Slepian--Wolf coding, syndrome coding, positional coding,
enumerative coding, entropy‑field coding, tensor coding, Minimum
Description Length, bits‑back coding, RNR Coding.

\textbf{1. Introduction}

\textbf{1.1 Motivation}

Classical lossless compressors --- LZ77/LZ78 derivatives, LZMA, the
Burrows--Wheeler family, arithmetic‑coded context mixers such as PAQ and
Cmix, and modern entropy coders based on arithmetic coding or asymmetric
numeral systems --- are highly effective on a wide range of inputs
because they exploit repeated substrings, local contexts, block sorting,
and adaptive probability models. Their operational question is
sequential: given the bytes seen so far, what is the next byte, or where
did this fragment occur previously? This question is well‑matched to
data whose redundancy is dominated by local repetition or short‑range
statistics.

Many practically important data classes are not well described by this
question. Database pages, raw or lightly preprocessed image sensor data,
telemetry grids, structured binary sections, scientific tensors, and
certain instruction streams contain structure that is positional,
geometric, multi-scale, or coupled across non-adjacent positions. For
such data, the conjecture motivating this paper is that sequential
redundancy is weaker than non-local structural regularity, and
dictionary-based methods leave entropy unexploited that more structural
predictors could in principle reach. The conjecture is empirical, not
proven here: §11 lists falsifiable bit-per-byte targets for specific
corpora, and the experimental-design companion document specifies the
pre-registered measurement protocol. Within this paper we make no
quantitative claim about absolute compression ratios on any specific
corpus.

Neural lossless compressors such as Bellard's NNCP, PixelCNN‑style
autoregressive models, integer discrete flows, and bits‑back schemes
(BB‑ANS, Bit‑Swap, HiLLoC) demonstrate that learned predictors can drive
arithmetic or numeral‑system coders to bit‑rates below those of strong
classical baselines on certain corpora. These methods almost universally
adopt the same operational question as classical coders --- predict the
next symbol --- and replace the predictor with a neural network.

A second obstruction blocks neural compressors from one of the largest
practical deployment niches for lossless coding: archive and
time-series storage engines inside database servers. Compressed
read-only filesystems (squashfs, EROFS), columnar archive engines, and
time-partitioned metric stores rely on \emph{random access} ---
reading a byte at offset $p$ in time polylogarithmic in archive
size $|X|$. Dictionary-based seekable formats (bgzip, zstd seekable)
achieve this operationally via independent compressed frames plus a
seek index. Autoregressive neural byte models (NNCP and similar) do
not, because their predictors require sequential evaluation of all
bytes $[0, p)$ before byte $p$ is available; the latent-variable lossless schemes
(IDF, Bit-Swap, HiLLoC) decode in block units that can be larger than
the desired random-access granularity. The author's engineering
experience with a Rust MySQL-compatible server motivated this
obstruction concretely (compressed-archive and time-series engines
across 20 storage engines): when ratio matters and seek latency
matters jointly, neither existing class fits the niche. This anecdotal
motivation is what suggested treating the neural model output as side
information rather than as a sequential predictor; the question
addressed in this paper is whether such a recasting can support random
access at sync-point granularity without losing the ratio advantage
(an empirical question deferred to §11).

This work proposes a different organizing question. Treat the neural
model output as side information about the source, not as a candidate
reconstruction. A neural generator alone cannot losslessly compress
data, because any single bit of disagreement between generator output
and ground truth is fatal. But if the model output is treated as a noisy
or partial description of the data --- a known damaged version, a
positional probability field, or a tensor of predicted symbol
probabilities and reliability --- then the archive needs to store only
the repair information that, combined with the model output, recovers
the exact original.

The Ratushnyi Neural Repair Coding (RNR) framework formalizes this idea
along a single axis: the dimensionality and target of prediction. Type‑I
predicts a one‑dimensional stream; Type‑II predicts a two‑dimensional
symbol‑position partition; Type‑III predicts a multi‑dimensional tensor
that may include the predicted entropy of prediction itself. The central
claim is conditional: for data classes in which neural side information
reduces residual uncertainty by more than the cost of describing the
model, the repair, the verification, and the metadata, RNR can in
principle outperform sequential dictionary or entropy‑only baselines.

\textbf{1.2 Position relative to existing work}

RNR draws on five established lines of work. Slepian and Wolf
established that a source can be encoded losslessly at a rate determined
by its conditional entropy given correlated side information at the
decoder. Syndrome‑based and parity‑check constructions, including LDPC
and BCH codes, give practical machinery for treating prediction error as
a virtual channel noise process. Cover's enumerative source coding shows
how subsets and sequences under combinatorial constraints can be encoded
by their rank among admissible possibilities. Rissanen's Minimum
Description Length principle provides the criterion under which a
learned representation justifies its own overhead. Arithmetic coding and
asymmetric numeral systems provide near‑optimal entropy coders for given
probability models.

Within neural lossless compression, the closest reference points are
NNCP (autoregressive byte‑level transformer with arithmetic coding),
PixelCNN++ and PixelSNAIL (autoregressive factorized image models with
categorical likelihoods), the bits‑back family (BB‑ANS, Bit‑Swap,
HiLLoC, encoding latent variables of a variational model losslessly via
ANS), integer discrete flows (IDF, IDF++), and L3C for hierarchical
image priors. Among classical context‑mixing coders, PAQ and Cmix
represent the upper end of what mixture‑of‑experts entropy coding
achieves without neural networks.

RNR is not a replacement for any of these. It is a taxonomy and a set
of constructions that, in our reading, naturally organize them ---
though the mappings below are interpretive (not theorems) and are
included to fix the vocabulary of later sections. Per-position
arithmetic coding over a factorized predicted distribution is
recovered as a special case of Type-I formally (see the reduction
sketch in §4.5). PixelCNN++ over images is autoregressive over pixels
conditioned on prior pixels (categorical or discretized-mixture
likelihoods, depending on the head), and the relation to Type-II is
analogy not equivalence: Type-II as defined in §3.2 is partition-based
and does not subsume PixelCNN++'s pixel-wise softmax. Bits-back coding
on a continuous latent fits into Type-III-C as the limit where the
discrete dictionary is replaced by a variational posterior, with
exact cross-entropy plus KL accounting (Theorem 6.6). PAQ-style
switching context mixing can be viewed as a hand-designed
Type-III-A-like scheduler with a continuous-weight logistic mix
(rather than a discrete argmin), so the mapping is analogy rather
than instantiation.
The framework's intended contribution is to give a unified vocabulary
and a set of constructions for three regimes that the existing methods
treat ad hoc or not at all --- semi-non-factorized positional models,
predicted entropy fields, and structured token tensors. We do not
claim that existing methods are empirically weak in these regimes;
Theorem 5.4 specifically separates only factorized-marginal coders
from semi-non-factorized ones, and a true autoregressive coder can in
principle match the joint-coding rate on the same source. Whether the
RNR constructions outperform existing methods on natural data classes
is the empirical question of §11; the present paper provides only the
theoretical scaffolding for the comparison.

\textbf{1.3 Contributions}

This article makes the following contributions.

\begin{itemize}
\item
  A framework definition (Section 3) that decouples the choice of
  predictor (the model M) from the choice of representation target
  (stream, partition, tensor) and the choice of repair encoding (mask,
  residual, rank, selector).
\item
  Concrete constructions for each of the five types: Type-I with an
  explicit Code12 systematic shortened-Hamming-like byte-to-12-bit
  mapping (Section 4); semi-non-factorized Type-II positional coding
  with disjoint-partition and count constraints (Section 5); and
  Type-III-A/B/C with entropy-field, byte-tensor, and hierarchical
  $32 \times 32$ token-tensor constructions (Section 6).
\item
  Achievability and converse bounds for Type-I (Theorems 4.1, 4.2):
  the Type-I rate matches the conditional source-coding lower bound
  $H(X \mid Y)$ to first order in $N$.
\item
  Complexity boundary for Type-I codeword selection: the fixed
  Code12 selection at $n_0 = 256, m_0 = 12$ is a finite constant-size
  optimization (Remark 4.3a), while the varying-parameter
  learned-code labeling problem (alphabet size $n$ and codeword
  length $m$ as input) is NP-hard via reduction from Wagner and
  Corneil's (1990) tree-to-hypercube subgraph embedding (Theorem 4.3).
\item
  A corrected positional-advantage theorem (Theorem 5.3) with a
  multinomial baseline and an explicit escape-mode cost term,
  addressing two errors in earlier drafts that overstated the
  Type-II advantage.
\item
  An amortized separation theorem (Theorem 5.4) for
  semi-non-factorized Type-II via a specific linear-code
  construction, with the exact information-theoretic quantity
  governing factorized-versus-joint coding identified as
  multi-information $\mathrm{TC}(D_N)$ (Theorem 5.5; the gap may be
  zero on sources without joint structure). Whether
  $\mathrm{TC}(D_N) = \Theta(N)$ holds for specific natural-data
  classes is an open empirical question (Remark 5.6 sketches a
  manifold-hypothesis intuition but does not establish a rigorous
  entropy bound).
\item
  A status remark on Type-II support selection
  (Remark 5.7): the natural marginal-$Q$-prefix heuristic is not
  provably optimal in the semi-non-factorized regime --- expected cost
  depends on the full joint law of the partition sets, not just
  marginals --- and a concrete counterexample is included in the
  paper's verification scripts. Optimal fixed-size support selection
  for general predictors remains open.
\item
  An approximate-oracle mode-selection bound for Type-III-A (Theorem
  6.1) and a tight matching lower bound (Theorem 6.2) showing the
  $2\varepsilon$ factor cannot be improved without additional
  assumptions on the prediction error structure.
\item
  A mutual-information characterization of the Type-III-B byte-tensor
  advantage (Theorem 6.3) and an invariance result (Theorem 6.4)
  showing all rooted binary tree decompositions of the 256-byte
  alphabet achieve the same information-theoretic rate.
\item
  Four results on Type-III-C: the greedy MDL admission criterion
  (Theorem 6.5), the bits-back integration with an exact
  cross-entropy plus KL-divergence accounting (Theorem 6.6),
  NP-hardness of the restricted (constant-predictor, normal-form
  recursive dictionary, symbol-count objective) dictionary-learning
  problem (Theorem 6.7), and APX-hardness of the same restricted
  case inherited from the Smallest Grammar Problem (Lehman--Shelat
  2002; Charikar et al.\ 2005) with constant $\gamma_{\mathrm{SGP}}
  > 1$ (Theorem 6.8). We do not pin down the specific numerical
  value of $\gamma_{\mathrm{SGP}}$, and we do not claim hardness of
  the unrestricted joint problem; both are listed as open.
\item
  A no-universal-dominance result (Theorem 8.1) establishing the
  framework's claim as conditional, and a verification-discipline
  theorem (Theorem 9.1) ruling out cross-subblock error propagation
  in autoregressive modes.
\item
  Sufficient conditions for bit-exact decoder reproduction under
  fixed predictors (Theorem 10.1) and adaptive predictors (Theorem
  10.2 with Definition 10.1, the compatible-online-optimizer class),
  reducing reproducibility from a software-engineering question to a
  conformance specification. Published low-precision-training methods
  (integer SGD, signSGD, integer Adam, integer-quantized inference)
  are candidate components but their conformance against the full
  specification on any specific hardware target requires implementation
  and testing.
\item
  Random-access (seek-and-read) capability with cost
  $O(\log(|X|/K) \cdot \log L(R) + (K + k) \cdot \mathrm{cost}(M))$
  --- polylogarithmic in archive size $|X|$ once sync-point spacing
  $K$, read window $k$, predictor cost, and $\log L(R)$ are held
  fixed (or polylog in $|X|$) --- via Theorem 10.3 with Definitions
  10.2, 10.3 of $W$-bounded predictors and sync points.
  Theorem 10.4 covers non-$W$-bounded predictors via state snapshots.
  The sync-point compression overhead for natural data is estimated
  in §10.7.4; a tight worst-case bound in general is not proved. The
  comparison to existing neural compressors (NNCP-style autoregressive
  byte models, block-decoding latent-variable schemes) on either seek
  capability or ratio is empirical and is the subject of §11.
\item
  Block-device and read-only filesystem deployment as direct
  corollaries (Theorems 10.5 and 10.6 in §10.8): an OS-level driver
  can serve LBA-indexed read requests with cost polylogarithmic in
  archive size $|X|$ in the parameter regime where the sync-point
  spacing $K$, read window $k$, predictor cost $\mathrm{cost}(M)$,
  and seek-index entry width $\log L(R)$ are fixed or grow at most
  polylogarithmically in $|X|$, and additionally
  $L(R) = \mathrm{poly}(|X|)$ (Theorem 10.5); and the predictor can
  be required to
  run bit-exactly on any hardware target satisfying the conformance
  specification of Theorem 10.1 (Theorem 10.6). Specific target
  classes named in the spec (CPU scalar, CPU SIMD including
  AVX-512, GPU integer Tensor Cores, integer NPUs) are engineering
  assertions about what targets satisfy the spec, not theorems;
  conformance for each named target requires implementation and
  testing per the conformance suite (engineering-spec document).
\item
  A runtime parameter-selection framework (Section 7.2) that subsumes
  input-dependent design choices (Code-$k$ rate, alphabet size,
  support threshold, tile size, dictionary size) under the existing
  mode-selection machinery, with the same $2\varepsilon$
  approximate-oracle gap as Theorem 6.1 plus the per-tile metadata
  overhead $O(\log K_{\text{params}})$ for transmitting the chosen
  parameter values.
\item
  Worked numerical examples for Type-I (§4.7), Type-II (§5.7), and
  Type-III-B (§6.4), demonstrating concrete coding-length arithmetic
  on toy distributions. Type-III-A and Type-III-C are illustrated
  through their composition and dictionary-construction sections
  respectively.
\end{itemize}

\textbf{1.4 Organization}

Section 2 fixes notation and reviews the necessary background. Section 3
defines the framework. Sections 4, 5, and 6 develop Type-I, Type-II,
and the three Type-III subtypes. Worked numerical examples are given
in §4.7 (Type-I), §5.7 (Type-II), and §6.4 (Type-III-B); Type-III-A
and Type-III-C are illustrated through the composition discussion in
§7 and the dictionary-construction discussion in §6.5 respectively
rather than through standalone worked examples. Section 7 discusses
composition, mode selection, and runtime parameter selection across
types. Section 8 states the
no‑universal‑dominance result. Section 9 covers verification and error
containment. Section 10 addresses practical reproducibility and proves a
sufficient condition for bit-exact decoder reproduction. Section 11
lists experimental predictions for empirical follow‑up. Section 12
surveys related work, and Section 13 reports the status of open
problems, distinguishing resolved results, remaining mathematical open
questions, and items explicitly deferred to engineering or empirical
work.

\textbf{2. Preliminaries and notation}

\textbf{2.1 The lossless compression problem}

A lossless compressor is an injective map
$\mathsf{Enc}$ from a source alphabet (here, finite-length binary
strings or byte sequences of a fixed block length $N$) to binary
codewords, together with a decoder $\mathsf{Dec}$ such that
$\mathsf{Dec}(\mathsf{Enc}(X)) = X$ for every $X$. Injectivity is
necessary; for the multi-stream concatenations used later (model,
repair, verification, metadata), each stream must in addition be
prefix-free or carry an explicit length so that the decoder can
partition the codeword.

Throughout the paper, we fix the block length $N$ and treat $X$ as a
random variable over the finite alphabet $\{0,1\}^{8N}$ (byte block)
under a distribution $D$ on that finite set. We write $L(\cdot)$ for
the bit length of a binary description, and adopt the convention
that $0 \cdot \log 0 = 0$. For a random variable $X$ over a finite
alphabet, $H(X)$ denotes the Shannon entropy of $X$ in bits. Given
a side-information variable $Y$, $H(X \mid Y)$ denotes the
conditional entropy.

For a source $X \sim D$ and a baseline compressor $B$, we say an RNR
coder $\mathsf{Enc}$ is conditionally advantageous on $D$ if
\[
\mathbb{E}[L(\mathsf{Enc}(X))] < \mathbb{E}[L(B(X))], \quad X \sim D.
\tag{2.1}
\]
\emph{Notation note.} The symbol $C$ in later sections denotes
\emph{public context} (the input to the predictor $M$, so
$Y = M(C)$), not the compressor; the symbol $D$ in later sections
denotes the source distribution, not the decoder. The
encoder/decoder pair is referred to as $\mathsf{Enc}/\mathsf{Dec}$
throughout to avoid this collision.

\textbf{2.2 Side-information coding}

Slepian and Wolf's distributed-source-coding theorem (Slepian and
Wolf 1973) characterizes the rate region for encoding jointly
distributed sources $X$ and $Y$ with separate encoders and a joint
decoder, with constraints $R_X \geq H(X \mid Y)$,
$R_Y \geq H(Y \mid X)$, and $R_X + R_Y \geq H(X, Y)$ in the
finite-alphabet memoryless setting. The corner of this region we use
is the asymmetric case in which $Y$ is already available to the
decoder, so only $X$ needs to be coded at rate
$\geq H(X \mid Y)$. The motivating intuition --- that decoder-side
information drives down the rate to a conditional entropy --- is
Slepian-Wolf's; the achievability in our setting, however, reduces
to ordinary \emph{conditional source coding} (Cover and Thomas 2006,
Theorem 5.4.1) and does not require the binning constructions that
Slepian-Wolf invokes when the encoder lacks $Y$.

Our setting differs from the classical distributed Slepian-Wolf
setting in three ways. First, we consider a single sender, not two
correlated sources at different physical locations. Second, the side
information $Y$ is not a separately observed correlated random
variable but the deterministic output of a public model $M$ applied
to a public context $C$: $Y = M(C)$. Third, this side information is
available to both encoder and decoder, so the encoder can exploit it
directly; the relevant achievability is conditional source coding,
and Slepian-Wolf binning is not needed.

We use the term ``side information'' in this broader sense
throughout: any deterministic, decoder-reproducible function of
public state that constrains the conditional distribution of $X$.

\textbf{2.3 Entropy coding}

We assume the availability of an entropy coder --- arithmetic coding
(Witten, Neal, Cleary 1987; Rissanen and Langdon 1979) or asymmetric
numeral systems (Duda 2009) --- that, given an explicit probability
model $Q$ with $Q(x) > 0$ on the support of the source and finite-
precision arithmetic of $b$ bits per state, encodes a sequence $x$ of
length $n$ in $\lceil -\log_2 Q(x) \rceil + \varepsilon_n$ bits,
where the redundancy $\varepsilon_n$ accounts for end-of-stream and
finite-precision rounding (under-2-bit total overhead for the
classical Witten-Neal-Cleary integer arithmetic coder; ANS adds an
$O(\log b)$ state-initialization overhead amortized over the stream).
Taking expectation under the true source $X \sim D$ gives expected
length $\mathbb{E}_D[-\log_2 Q(X)] + \varepsilon_n$, i.e.,
cross-entropy plus redundancy. We assume $\varepsilon_n / n \to 0$
as the block length $n$ grows.

\textbf{2.4 Enumerative coding}

Cover's enumerative source coding (Cover 1973) describes an element
of a \emph{known} finite class by its rank among the elements of that
class. The information-theoretic cost is the binary logarithm of the
class size; a self-delimiting binary rank requires
$\lceil \log_2 |\mathcal{C}| \rceil$ bits (or
$\log_2 |\mathcal{C}| + O(1)$ via arithmetic coding driven by the
uniform rank distribution). Whenever we write the cost as
$\log_2 |\mathcal{C}|$ below we mean the information-theoretic rate;
the integer-representation overhead is absorbed into the
$\varepsilon_n$ entropy-coder redundancy of §2.3. Crucially, the
class itself must be decoder-reproducible: when its description
depends on data-dependent parameters (e.g., the cardinality vector
of a partition), those parameters must be transmitted separately.
Two enumerative formulas are used repeatedly in this paper.

The number of subsets of size $k$ drawn from $\{1, \ldots, N\}$ is
$\binom{N}{k}$, so a subset is described in
$\log_2 \binom{N}{k}$ bits relative to the class
$\binom{[N]}{k}$ \emph{given $k$}. When $k$ is data-dependent it is
transmitted in $O(\log N)$ bits before the rank.

The number of partitions of $\{1, \ldots, N\}$ into 16 disjoint sets
$M_0, \ldots, M_{15}$ with cardinalities $k_0, \ldots, k_{15}$ is the
multinomial coefficient
\[
M(N; k_0, \ldots, k_{15}) = \frac{N!}{k_0!\, k_1! \cdots k_{15}!},
\tag{2.2}
\]
so such a partition is described in
$\log_2 M(N; k_0, \ldots, k_{15})$ bits given the count vector
$(k_0, \ldots, k_{15})$. The count vector itself lies in a set of
size at most $(N + 1)^{16}$ and is transmitted in
$16 \log_2(N + 1) = O(\log N)$ bits by direct enumeration. The
multinomial baseline is strictly tighter than
$\sum_v \log_2 \binom{N}{k_v}$, which corresponds to encoding 16
independent subsets without enforcing disjointness, and therefore
overcounts.

\textbf{2.5 Minimum Description Length}

Rissanen's Minimum Description Length principle (Rissanen 1978)
selects the model that minimizes the total description length: model
structure, real-valued parameters, and data given the model. There
are several technical formulations --- crude two-part MDL,
prequential MDL, universal coding, normalized maximum likelihood
(NML) --- with different optimality criteria. RNR uses MDL only
heuristically, in its simplest two-part admission form: a candidate
model or token is included if its description cost plus the resulting
conditional code length beats the alternative without it. The
admission rule
\[
L(M) + L(X \mid M) < L_{\text{base}}(X)
\]
is a two-part-MDL comparison heuristic, not a theorem of optimality;
we invoke it as a design criterion at every level (Type-I model
admission, Type-III-C token admission) but do not claim that any RNR
construction is the MDL optimum under a stronger criterion such as
NML.

\textbf{2.6 Bits-back coding}

Bits-back coding is a technique for losslessly compressing data under
a latent-variable model. The encoder uses an auxiliary bit string to
sample a latent $z$ from a variational posterior $q(z \mid x)$,
encodes $(z, x)$ jointly via prior $P(z)$ and conditional
$P(x \mid z)$, and the decoder recovers $z$ and $x$ and extracts the
auxiliary bits. The expected net length under the data distribution
$X \sim D$ and the variational sampler $z \sim q(\cdot \mid X)$ is
the \emph{negative ELBO}:
\[
\mathbb{E}_{X \sim D,\, z \sim q(\cdot \mid X)} \!\left[-\log P(X, z) + \log q(z \mid X)\right]
= \mathbb{E}_{X \sim D}\!\left[-\log P_X(X)\right]
+ \mathbb{E}_{X \sim D}\!\left[\mathrm{KL}(q(\cdot \mid X) \Vert P(\cdot \mid X))\right].
\]
This equals the source entropy $H_D(X)$ only when the model marginal
matches the data ($P_X = D$) and $q$ equals the true posterior; in
general it is cross-entropy under the model marginal plus the
variational gap. Theorem 6.6 makes this accounting explicit for the
RNR Type-III-C decomposition. The lineage is: Hinton and van Camp
(1993) introduced the bits-back argument for noisy-weight
communication in a Bayesian-network MDL setting; Frey (1997, 1998)
extended it to one-to-many stochastic source codes with discrete
latents in graphical models; Townsend, Bird and Barber (2019) gave
the practical BB-ANS construction that handles continuous latents via
ANS coding (Bit-Swap and HiLLoC are subsequent extensions). We use
the term ``bits-back'' for the general principle and BB-ANS for the
specific continuous-latent realization; see Section 6.5.

\textbf{2.7 Notation summary}

We adopt the following notation throughout.
\begin{description}
\item[$N$] block length (number of bytes per source block, fixed throughout)
\item[$n$] generic symbol-sequence length (e.g., positional, when distinct from $N$)
\item[$X$] original source block (sequence of bytes, nibbles, or symbols)
\item[$D$] source distribution on $\{0,1\}^{8N}$; $X \sim D$
\item[$Y$] decoder-reproducible side information: $Y = M(C)$
\item[$M$] deterministic predictor; $M_I, M_{\text{pos}}$ for type-specific predictors
\item[$C$] public context (seed, verified previous subblock, metadata), distinct from the §2.1 compressor symbol
\item[$Q$] predicted probability field (vector, matrix, or tensor); used as the entropy-coder model in §2.3
\item[$U_m(t)$] predicted code-length / mode-cost field for tile $t$ under mode $m$ (Type-III-A)
\item[$T$] invertible representation transform; $Z = T(X)$
\item[$R$] repair / residual / rank / selector stream
\item[$V$] verification overhead (hash, checksum); $H_v(X)$ stored as $v$
\item[$\Theta$] metadata overhead (sizes, mode flags, quantization tables)
\item[$B$] baseline compressor (LZMA, zstd, NNCP, $\ldots$); $L_{\text{base}}$ is its description length
\item[$\mathsf{Enc}, \mathsf{Dec}$] encoder/decoder pair (compressor); see §2.1
\item[$\varepsilon_n$] entropy-coder redundancy over $n$ symbols (positivity + finite-precision overhead per §2.3)
\item[$\Pi(X) = (M_0, \ldots, M_{15})$] positional partition of $X$ into symbol-class position sets
\item[$k_v$] cardinality of $M_v$, i.e., count of symbol $v$ in the block
\item[$H(\cdot), H(\cdot \mid \cdot)$] Shannon entropy / conditional entropy in bits
\item[$I(\cdot\,;\cdot \mid \cdot)$] conditional mutual information
\end{description}

\textbf{3. The RNR framework}

\textbf{3.1 The repair principle}

A neural generator alone cannot losslessly compress data, because the
generator output and the source disagree on at least some bits with
non‑negligible probability, and a single bit of disagreement is fatal.
RNR resolves this by treating the model output not as a candidate
reconstruction but as side information, and storing a separate repair
description that combined with the side information reconstructs the
exact source. The framework is parameterized by three independent
choices.

First, the predictor M. This may be a small autoregressive transformer,
a positional convolutional model, a learned token predictor, a
non‑neural context‑mixing predictor (PAQ‑style), a hash‑based predictor,
or any deterministic function of public state.

Second, the representation target. The model may predict (a) a
one‑dimensional stream of bytes or codewords, (b) a two‑dimensional
positional probability field over a partition of positions among symbol
classes, or (c) a higher‑dimensional tensor that may include the model's
own predicted uncertainty as a field. These three choices define Type‑I,
Type‑II, and Type‑III respectively.

Third, the repair encoding. Given a representation target and a model
prediction, the residual that the archive stores may be (i) an error
mask between predicted and true codewords, (ii) a rank within a
constrained enumerative class, (iii) a selector from a candidate list,
(iv) a per‑position probability under which an entropy coder operates,
(v) a tile‑wise mode flag, or any composition of these. The repair
encoding is a design choice within each type.

\textbf{3.2 The five RNR types}

\textbf{Definition 3.1 (RNR Type-I).} An RNR Type-I coder predicts a
one-dimensional damaged stream $Y = M_I(C)$ and stores a repair
stream $R_I$ together with a verification value $v = H_v(X)$ as part
of the archive. The decoder forms
$\hat{X} = \mathrm{DecRepair}(Y, R_I)$ and accepts the reconstruction
only if $H_v(\hat{X}) = v$; for a valid archive, $\hat{X} = X$
exactly. The canonical instance is Code12, in which each source byte
is mapped to a 12-bit codeword and the predicted error mask is
entropy-coded by Hamming-weight class.

\textbf{Definition 3.2 (RNR Type-II).} An RNR Type-II coder predicts
a positional marginal probability field $Q_v(i)$ approximating
$\Pr(i \in M_v \mid C)$, $v \in \{0, \ldots, 15\}$,
$i \in \{1, \ldots, N\}$, with $\sum_v Q_v(i) = 1$ for each position
$i$. The archive stores a partition residual $R_\Pi$ that, combined
with $Q$, recovers the actual partition
$\Pi(X) = (M_0, \ldots, M_{15})$, where the sets $M_v$ are pairwise
disjoint subsets of $\{1, \ldots, N\}$ and
$\sum_{v=0}^{15} |M_v| = N$. The recovered partition determines the
source $X$.

\textbf{Definition 3.3 (RNR Type-III-A).} An RNR Type-III-A coder
predicts an entropy field $U_m(t)$ over tiles $t$ and coding modes
$m$, selects predicted per-tile modes
\[
\hat{m}(t) = \min \{ m : U_m(t) = \min_{m'} U_{m'}(t) \},
\]
breaking ties by lexicographic order on a fixed canonical mode
index (specified by the model description as a total ordering of
mode identifiers; e.g., the ordering in which modes are declared in
$L(M)$). For decoder reproducibility under finite-precision
arithmetic, the $U_m(t)$ values must be computed in the same
deterministic representation on encoder and decoder (e.g., as the
bit-exact integer outputs of the same conforming inference engine
per Theorem 10.6); two encoders that produce $U_m(t)$ in
floating-point and round differently would tie-break differently
and break decoder reproducibility. The selection is then
deterministic and decoder-reproducible. The archive stores the
per-tile residual streams and, when the encoder departs from the
predicted argmin, override flags as metadata.

\textbf{Definition 3.4 (RNR Type-III-B).} An RNR Type-III-B coder
predicts a normalized byte tensor $Q_{h, l}(i)$ approximating
$\Pr(h_i = h, l_i = l \mid C)$ over high-nibble $h$ and low-nibble $l$
of byte $i$, with $h, l \in \{0, \ldots, 15\}$ and
$\sum_{h, l} Q_{h, l}(i) = 1$ for each position $i$. A
decoder-reproducible mode choice, either stored as metadata or
computed by a fixed rule from $(C, Q)$, selects direct 256-class
partition coding or a hierarchical scheme that first encodes the
high-nibble partition and then the low-nibble partition conditional
on the high nibble.

\textbf{Definition 3.5 (RNR Type-III-C).} An RNR Type-III-C coder is
parameterized by a learned token dictionary
$D = \{\tau_0, \ldots, \tau_{K-1}\}$ with $K = a \cdot b$ (typically
$32 \times 32 = 1024$ or $64 \times 64 = 4096$) hierarchically indexed
by $\text{family} \in \{0, \ldots, a-1\}$ and
$\text{variant} \in \{0, \ldots, b-1\}$ as
$\text{token\_id} = b \cdot \text{family} + \text{variant}$. The
model predicts $Q(k, i)$ approximating
$\Pr(\tau_k \text{ begins at position } i \mid C)$.
The archive stores a tokenization residual $R_C$ sufficient to
reconstruct the exact token-interval cover of the block.

\textbf{3.3 The repair length decomposition}

Every RNR coder has the same general code-length decomposition:
\[
L_{\text{RNR}}(X) = L(M) + L(R \mid Y) + L(V) + L(\Theta),
\tag{3.1}
\]
where $L(M)$ is the model description length (parameters,
quantization tables, adapter weights as relevant), $L(R \mid Y)$ is
the repair stream length, $L(V)$ is verification overhead, and
$L(\Theta)$ is metadata overhead (sizes, mode flags, dictionary
headers). The conditional advantage criterion is then
\[
\mathbb{E}_D\!\left[L(M) + L(R \mid Y) + L(V) + L(\Theta)\right] < \mathbb{E}_D[B(X)].
\tag{3.2}
\]

\textbf{3.4 Amortization of model cost}

$L(M)$ is paid once and amortized across all blocks encoded by the
same model. For a single-block archive, $L(M)$ dominates and the
framework loses unless $M$ is extremely small. For an archive of $m$
blocks of average size $n$, the effective per-block model cost is
$L(M)/m$, and the criterion (3.2) becomes more permissive in $m$.
Throughout, we assume the amortized regime; Section 10 discusses the
single-block case explicitly.

\textbf{4. Type-I: Stream Repair Coding}

\textbf{4.1 Definition and operational form}

Type-I is the one-dimensional form. The source
$X = (x_1, \ldots, x_N)$ is a stream of bytes or symbols. The model
produces a predicted stream $Y = M_I(C)$ of the same length,
considered as a damaged version of $X$. The encoder, knowing both
$X$ and $Y$, computes a repair $R_I = \mathrm{EncRepair}(X, Y)$, and
the decoder recovers $\hat{X} = \mathrm{DecRepair}(Y, R_I)$. The
encoder also stores the verification value $v = H_v(X)$ in the
archive; the decoder accepts the reconstruction only if
$H_v(\hat{X}) = v$ (Definition 3.1).

\textbf{4.2 The Code12 construction}

In the Code12 variant of Type-I, each source byte
$b \in \{0, \ldots, 255\}$ is mapped through an injective function
$E_{12} : \{0, \ldots, 255\} \to \{0,1\}^{12}$ to a 12-bit codeword
$c = E_{12}(b)$. The model predicts a codeword $\hat{c}_i = M_I(C, i)$
for each position. The raw error mask is
$e_i = c_i \oplus \hat{c}_i$. The motivation is that, with a
well-chosen $E_{12}$, the error-mask distribution concentrates on
low Hamming weights when the model is approximately correct, and
Hamming weight is cheap to entropy-code.

We specify $E_{12}$ as a shortened-Hamming-like systematic mapping. Let
the source byte bits be $(b_0, \ldots, b_7)$. Use 1-based codeword
positions $c_1, \ldots, c_{12}$ and place data bits at non-parity
positions:
\[
\begin{aligned}
c_3 &= b_0, & c_5 &= b_1, & c_6 &= b_2, & c_7 &= b_3, \\
c_9 &= b_4, & c_{10} &= b_5, & c_{11} &= b_6, & c_{12} &= b_7.
\end{aligned}
\]
The parity positions $c_1, c_2, c_4, c_8$ are computed by even parity
over selected data positions:
\[
\begin{aligned}
c_1 &= c_3 \oplus c_5 \oplus c_7 \oplus c_9 \oplus c_{11}, \\
c_2 &= c_3 \oplus c_6 \oplus c_7 \oplus c_{10} \oplus c_{11}, \\
c_4 &= c_5 \oplus c_6 \oplus c_7 \oplus c_{12}, \\
c_8 &= c_9 \oplus c_{10} \oplus c_{11} \oplus c_{12}.
\end{aligned}
\]
This construction is reversible because the data bits $b_0, \ldots, b_7$
appear directly in $c$. We use it for two reasons: first, single-bit
and many two-bit codeword errors decode to small Hamming-weight masks,
which are cheap; second, the parity structure means random predictor
noise distributes errors over both data and parity positions, rather
than concentrating in the data bits. We discuss alternative $E_{12}$
mappings (Gray-ordered, learned, bucketed with candidate lists) in
Section 4.6.

\textbf{4.3 Repair stream encoding}

The simplest repair symbol at position $i$ is the error mask $e_i \in
\{0,1\}^{12}$ itself, of which there are $2^{12} = 4096$ possibilities.
To avoid paying 12 bits per position regardless of prediction quality,
we encode by Hamming-weight class:
\begin{itemize}
\item class 0: no error (1 codeword);
\item class 1: one-bit flip (12 codewords, $\log_2 12 \approx 3.585$ bits);
\item class 2: two-bit pattern ($\binom{12}{2} = 66$ codewords, $\log_2 66 \approx 6.044$ bits);
\item class candidate-$k$: correct codeword is $k$-th in predicted candidate list;
\item class escape: store raw byte (8 bits) or raw codeword (12 bits).
\end{itemize}
Let the class probabilities be $p_0, p_1, p_2, p_{\text{cand}}, p_{\text{esc}}$
summing to 1. The expected per-position repair length is

\[
\ell_{\text{rep}} = p_0 \cdot 0 + p_1 \log_2(12) + p_2 \log_2(66) + p_{\text{cand}} \cdot (\text{header} + \text{rank}) + p_{\text{esc}} \cdot 8 + H(\text{class}) + \varepsilon,
\tag{4.1}
\]

where $H(\text{class})$ is the entropy of the class label encoded with
a small adaptive entropy coder, and $\varepsilon$ is the entropy-coder
redundancy. When $p_0$ dominates, the dominant cost is $H(\text{class})$
plus a tiny contribution from the other classes; $\ell_{\text{rep}}$
can fall well below 1 bit per byte if the model is good.

\textbf{4.4 Achievability bound}

\textbf{Theorem 4.1 (Type-I achievability).} \emph{Let $X$ be a source
block of $N$ bytes, $Y = M_I(C)$ be the predicted Code12 stream, and
$E = (e_1, \ldots, e_N)$ be the error-mask sequence with empirical
class distribution $P_E$. Suppose the class entropy coder achieves
expected length at most $N \cdot H(P_E) + \varepsilon_N$. Then the
expected Type-I code length is bounded by}
\[
\mathbb{E}[L_I(X)] \leq L(M_I) + N \cdot H(P_E) + N \cdot \ell_{\text{cond}} + \varepsilon_N + L(V) + L(\Theta),
\tag{4.2}
\]
\emph{where $\ell_{\text{cond}}$ denotes the expected conditional
payload (within-class rank or escape) per position. Consequently,
Type-I is conditionally advantageous on a distribution $D$ if the
right-hand side, in expectation under $D$, is smaller than the
expected baseline length $\mathbb{E}_D[B(X)]$.}

\emph{\textbf{Proof.}} The Type-I code is the concatenation of the
model description, the class stream, the within-class payloads, the
verification block, and the metadata block. By the entropy-coder
hypothesis, the class stream has expected length at most
$N \cdot H(P_E) + \varepsilon_N$. The within-class payload is a
separate stream whose expected length, summed across positions, is
$N \cdot \ell_{\text{cond}}$ by definition. Adding the constant
overheads $L(M_I), L(V), L(\Theta)$ gives the bound. Taking
expectation under $D$ and comparing to $\mathbb{E}_D[B(X)]$ yields
the advantage condition. \hfill$\blacksquare$

\textbf{Theorem 4.2 (Type-I converse and ideal achievability).}
\emph{Let $X \sim D$ be a source over length-$N$ byte blocks and let
$Y = M_I(C)$ be decoder-reproducible side information. Fix an
\emph{injective} Code12 map
$E_{12} : \{0, \ldots, 255\} \to \{0, 1\}^{12}$ so that, given $Y$,
the error variable $E = E_{12}(X) \oplus Y$ stands in bijection with
$X$ (the decoder recovers $X$ from $(Y, E)$ via
$X = E_{12}^{-1}(E \oplus Y)$). Then any uniquely decodable lossless
code for $X$ whose decoder uses only the codeword and $Y$ has expected
length at least $H(X \mid Y)$. Conversely, an \emph{ideal} Type-I
class-and-payload coder driven by the true conditional law of $E$
achieves expected length $H(X \mid Y) + \varepsilon_N$, plus the
constant headers $L(M_I) + L(V) + L(\Theta)$. A concrete coder using
an approximate class model or a uniform within-class payload pays the
corresponding cross-entropy, which matches the bound only when those
models match the true conditional law. Non-injective $E_{12}$ would
collapse the bijection: e.g., if $E_{12}(0) = E_{12}(1) = 0$ and
$Y \equiv 0$, then $E \equiv 0$ for $X \sim \mathrm{Unif}\{0, 1\}$
while $H(X \mid Y) = 1$, so no decoder can recover $X$ from $(Y, E)$.}

\emph{\textbf{Proof.}} The converse is obtained by applying Shannon's
source-coding bound to each conditional distribution
$P(X \mid Y = y)$ and averaging over $Y$ (Cover and Thomas 2006,
Theorem 5.4.1; see also Chapter 15 for coding with side information).
For achievability, fix an injective Code12 map and define
$E = E_{12}(X) \oplus Y$. Given $Y$, the map $X \leftrightarrow E$ is
bijective, so $H(E \mid Y) = H(X \mid Y)$. Encoding first the class
$C_E$ of $E$ and then the payload within that class, both under their
true conditional probabilities given $Y$, achieves
$H(C_E \mid Y) + H(E \mid C_E, Y) = H(E \mid Y) = H(X \mid Y)$ by the
chain rule for entropy (Cover and Thomas 2006, Theorem 2.6.4). An entropy coder driven by these conditional laws attains
this rate up to arithmetic-coding redundancy $\varepsilon_N$. A
specific Code12 implementation that uses an approximate $P(C_E \mid Y)$
or a uniform within-class payload (the construction of Theorem 4.1)
pays the corresponding cross-entropy in place of the entropy; it
equals the bound only when those models match the true conditional
law. The constant headers $L(M_I) + L(V) + L(\Theta)$ are $O(1)$
independent of $N$. \hfill$\blacksquare$

\textbf{4.5 Reduction to NNCP}

Consider the degenerate Code12 mapping in which $c$ is simply the byte
$b$ zero-padded to 12 bits, and the model predicts a probability
distribution $Q_i$ over the next byte rather than a single damaged
codeword. Then the optimal Code12 repair is to discard the predicted
codeword and encode each position via arithmetic coding under $Q_i$,
achieving expected per-position cost equal to the categorical
cross-entropy plus $\varepsilon/N$. This recovers NNCP exactly. Type-I
therefore generalizes NNCP by allowing (a) non-trivial $E_{12}$
mappings whose error distribution is more concentrated than a flat
softmax, and (b) class-based entropy coding that can be cheaper than
arithmetic coding for sharply peaked predictors.

\textbf{4.6 Variants of $E_{12}$}

Four $E_{12}$ variants are worth comparing empirically.
\begin{itemize}
\item
  $E_{12}$-A: shortened Hamming(12,8)-like systematic mapping as in
  Section 4.2.
\item
  $E_{12}$-B: Gray-ordered mapping in which numerically adjacent bytes
  differ by Hamming weight 1, suitable when model errors are
  predominantly small numerical perturbations of the true byte.
\item
  $E_{12}$-C: learned mapping in which the byte-to-codeword assignment
  is optimized to minimize the expected Hamming weight of $c \oplus
  \hat{c}$ over a training corpus and predictor.
\item
  $E_{12}$-D: bucketed mapping in which bytes that the predictor
  frequently confuses share a 4-byte bucket, and the repair stores a
  2-bit within-bucket selector.
\end{itemize}

Among these, only $E_{12}$-A and $E_{12}$-B admit closed-form
construction. $E_{12}$-C and $E_{12}$-D are learned, and we now show
that finding the optimum is computationally intractable.

\textbf{Remark 4.3a (Fixed-size Code12 search).} In the deployed
Code12 setting the alphabet size and code length are fixed constants:
$n_0 = 256$ and $m_0 = 12$. For these fixed parameters, optimizing
$E_{12}$ is a finite search over $\binom{2^{12}}{256} \cdot 256!$
injective byte-to-codeword maps, a number independent of the input
description length of the empirical confusion data. Exhaustive search
is therefore astronomically large but still a constant-size
enumeration, hence polynomial in the formal asymptotic sense. Thus
the practical Code12 case is not itself an asymptotically NP-hard
problem; NP-hardness only applies to a varying-parameter family in
which $n$ or $m$ is part of the input.

\textbf{Theorem 4.3 (NP-hardness of varying-parameter learned-code
labeling).} \emph{Consider the family of instances specified by
integers $n, m$ with $n \leq 2^m$, a symmetric row-stochastic matrix
$F \in \mathbb{Q}_{\geq 0}^{n \times n}$, and a rational threshold
$B$. Given an injective labeling
$E : \{1, \ldots, n\} \to \{0, 1\}^m$, define the unnormalized
uniform-prior Hamming loss}
\[
C_F(E) = \sum_{i,j=1}^n F_{ij} \cdot d_H(E(i), E(j)).
\]
\emph{Deciding whether there exists an injective labeling $E$ with
$C_F(E) \leq B$ is NP-hard in the description size of $(F, B, n, m)$.
Consequently, optimizing the learned code labeling is NP-hard for
varying parameters.}

\emph{\textbf{Proof.}} Reduce from tree-to-hypercube subgraph
embedding. Let $Q_m$ denote the graph on $\{0,1\}^m$ in which two
vertices are adjacent exactly when their Hamming distance is one.
Wagner and Corneil (1990) showed that deciding, given a tree $T$ and
integer $m$, whether $T$ is isomorphic to a subgraph of $Q_m$ is
NP-complete. Equivalently, the question is whether there is an
injective map $\phi : V(T) \to \{0, 1\}^m$ such that
$d_H(\phi(u), \phi(v)) = 1$ for every tree edge $\{u, v\}$.

If $|V(T)| > 2^m$, then no injective embedding into $Q_m$ exists. In
that case the reduction outputs the fixed no-instance with $n' = 2$,
$m' = 1$,
\[
F' = \begin{pmatrix} 0 & 1 \\ 1 & 0 \end{pmatrix}, \qquad B' = 1,
\]
which satisfies $n' \leq 2^{m'}$ and has no injective labeling with
$C_{F'}(E) \leq B'$ (the only injective $E : \{1, 2\} \to \{0, 1\}$
gives $C_{F'}(E) = 2 > B'$). Hence we may assume for the remaining
construction that $|V(T)| \leq 2^m$.

Given such an instance, write $n = |V(T)|$ and
$\Delta = \max_{v \in V(T)} \deg_T(v)$. For nontrivial instances
$\Delta \geq 1$. Construct an $n \times n$ matrix $F$ indexed by the
vertices of $T$ as follows:
\[
F_{uv} =
\begin{cases}
1/\Delta, & u \neq v \text{ and } \{u, v\} \in E(T), \\
1 - \deg_T(u)/\Delta, & u = v, \\
0, & \text{otherwise.}
\end{cases}
\]
This construction is polynomial in the size of $T$. The matrix is
symmetric because $T$ is undirected. It is row-stochastic because, for
each vertex $u$,
\[
F_{uu} + \sum_{v : \{u, v\} \in E(T)} F_{uv}
= 1 - \deg_T(u)/\Delta + \deg_T(u)/\Delta = 1.
\]
Thus $F$ is a valid symmetric row-stochastic confusion matrix: the
off-diagonal mass follows the tree edges, and the remaining
probability is placed on the diagonal self-confusion term.

Set
\[
B = \frac{2(n - 1)}{\Delta}.
\]
For any injective labeling $E$, diagonal terms contribute zero
($d_H(E(u), E(u)) = 0$), so
\[
C_F(E) = \frac{1}{\Delta}
\sum_{\{u, v\} \in E(T)} \!\!\bigl(d_H(E(u), E(v)) + d_H(E(v), E(u))\bigr)
= \frac{2}{\Delta} \sum_{\{u, v\} \in E(T)} d_H(E(u), E(v)).
\]
Since $E$ is injective, every tree edge has distinct endpoint labels,
hence $d_H(E(u), E(v)) \geq 1$. Because a tree on $n$ vertices has
$n - 1$ edges, every feasible labeling satisfies
\[
C_F(E) \geq \frac{2(n - 1)}{\Delta} = B.
\]
Equality holds if and only if every tree edge is mapped to Hamming
distance one. Therefore $T$ embeds as a subgraph of $Q_m$ if and only
if the constructed learned-code instance has an injective labeling
with $C_F(E) \leq B$. A polynomial-time algorithm for the learned-code
decision problem, or for exact optimization followed by comparison
with $B$, would decide the NP-complete tree-to-hypercube embedding
problem of Wagner and Corneil (1990). Hence varying-parameter
learned-code labeling is NP-hard. \hfill$\blacksquare$

\textbf{Remark 4.3b (Handling asymmetric confusion matrices).} For any
nonnegative confusion matrix $F$,
\[
\sum_{i, j} F_{ij} \cdot d_H(E(i), E(j))
= \sum_{i, j} \frac{F_{ij} + F_{ji}}{2} \cdot d_H(E(i), E(j)),
\]
because $d_H$ is symmetric. Thus an asymmetric empirical confusion
matrix induces the same labeling objective as its symmetric pair-weight
matrix $(F + F^T)/2$. However, if $F$ is row-stochastic,
$(F + F^T)/2$ need not be row-stochastic, so this identity does not by
itself extend the NP-hardness statement of Theorem 4.3 to the
asymmetric row-stochastic promise.

Consequently, $E_{12}$-C and $E_{12}$-D must be obtained by
approximation algorithms. Standard relaxations include gradient
descent on a continuous relaxation of the assignment (interpreting
$E_{12}$ as a stochastic mapping with simplex constraints), simulated
annealing, and tabu search. These approximations carry no
constant-factor guarantee in the varying-parameter regime; on the
fixed Code12 instance they are evaluated empirically against the
constant-size oracle of Remark 4.3a. The paper makes no
inapproximability claim for $E_{12}$: a sound constant-gap reduction
into the Hamming-restricted geometry is left as an open problem.

\textbf{4.7 Worked numerical example}

Consider a block of $N = 4096$ bytes drawn from a source where the
Code12 predictor achieves the following empirical class distribution:
$p_0 = 0.60$ (no error), $p_1 = 0.25$ (one-bit flip), $p_2 = 0.10$
(two-bit pattern), $p_{\text{cand}} = 0.04$ (within candidate list of
size 4), $p_{\text{esc}} = 0.01$ (escape to raw byte).

The class entropy is
\[
H(P_E) = -0.60 \log_2 0.60 - 0.25 \log_2 0.25 - 0.10 \log_2 0.10
- 0.04 \log_2 0.04 - 0.01 \log_2 0.01 \approx 1.527 \text{ bits.}
\]
The expected within-class payload per position is
\[
\ell_{\text{cond}} = 0.25 \log_2 12 + 0.10 \log_2 66 + 0.04 \log_2 4
+ 0.01 \cdot 8 \approx 0.896 + 0.604 + 0.080 + 0.080 = 1.660 \text{ bits.}
\]
Per-position total is therefore $1.527 + 1.660 \approx 3.19$ bits, or
$3.19 \cdot 4096 \approx 13063$ bits $= 1633$ bytes for the repair
stream. Adding $L(V) = 32$ bytes (a SHA-256 per 4 KiB subblock) and
$L(\Theta) = 16$ bytes of metadata gives 1681 bytes plus the amortized
model cost. For comparison, a 4096-byte block at 8 bits per byte raw
is 4096 bytes; LZMA on structured text typically achieves 1500--2000
bytes. The Type-I coder is competitive whenever the amortized $L(M_I)$
per block is below approximately 100--300 bytes, which is achievable
for a small distilled byte-level predictor.

\textbf{5. Type‑II: Positional Symbol Coding}

\textbf{5.1 From sequential to positional prediction}

Type‑I and the classical neural lossless coders share an operational
question: predict the next symbol. Type‑II changes the unit of modeling.
Instead of asking ``what is the next symbol?'', Type‑II asks: ``where
does each symbol appear in the block?'' The block is represented as a
partition of position indices among symbol classes, and the model
predicts the geometry of that partition.

This change is useful when a block has positional, geometric,
matrix‑like, periodic, or otherwise non‑local structure that is not
visible as local dictionary repetition. Examples include raw image data,
Bayer‑mosaic sensor data, database pages with structured record offsets,
instruction streams with positional opcode constraints, and numeric
matrices.

\textbf{5.2 Nibble partition representation}

Without loss of generality, we present Type-II at the nibble alphabet
$A_{16} = \{0, \ldots, 15\}$. Larger alphabets (full bytes, byte pairs,
learned tokens) generalize directly and are revisited in Section 6.
Let $X = (x_1, \ldots, x_N)$ with $x_i \in A_{16}$. For each nibble
value $v$, define the position set
\[
M_v = \{\, i \in \{1, \ldots, N\} : x_i = v \,\}.
\]
The collection $\Pi(X) = (M_0, \ldots, M_{15})$ is a partition of
$\{1, \ldots, N\}$: the $M_v$ are pairwise disjoint and their union is
the full position set. The original sequence $X$ is recovered from
$\Pi(X)$ by reading off the position memberships in order. Encoding
$X$ is therefore equivalent to encoding $\Pi(X)$.

\textbf{5.3 Factorized and semi-non-factorized positional models}

A Type-II model predicts a probability field $Q \in [0,1]^{16 \times N}$
with column normalization $\sum_v Q_v(i) = 1$. We distinguish two
regimes that affect the framework's relationship to Type-I.

In the factorized regime, $Q_v(i)$ is treated as if the position labels
were independent across $i$: the joint probability assigned to a
candidate sequence $(x_1, \ldots, x_N)$ is the product over $i$ of
$Q_{x_i}(i)$. Under this regime, optimal entropy coding achieves
length equal to $-\sum_i \log_2 Q_{x_i}(i)$, which is bit-equivalent
to per-position arithmetic coding driven by the same $Q$. Type-II
under a factorized model is therefore not operationally distinct from
Type-I with a per-position predictor. We make this explicit in
Lemma 5.1.

In the semi-non-factorized regime, the model additionally enforces
structural constraints that the factorized form violates. The minimal
such constraints are: (a) the disjoint-partition constraint that
exactly one $v$ has $x_i = v$ at each $i$; and (b) the count
constraint that the partition cardinalities $k_v = |M_v|$ are
themselves modeled (either predicted, transmitted, or constrained by a
learned distribution over count vectors). Under the
semi-non-factorized regime, the joint distribution over partitions is
not a product of per-position categoricals, and the coding gain over
the factorized regime can be substantial when the counts or the joint
geometry are highly informative. Theorem 5.4 quantifies this
separation.

\textbf{5.4 Coding strategies for the partition residual}

Given $Q$ and possible knowledge of counts $k_v$, the archive must
store a residual $R_\Pi$ sufficient to recover the exact partition.
Several coding strategies are available.
\begin{itemize}
\item
  Per-position arithmetic coding under $Q$: encode $x_i$ with
  probability $Q_{x_i}(i)$. Cost per position: $-\log_2 Q_{x_i}(i)$.
\item
  Per-symbol enumerative coding given the counts: for each $v$ in
  some order, encode $M_v$ as a subset of the remaining positions by
  its rank among $\binom{N'}{k_v}$ possibilities, where $N'$ is the
  number of positions not yet assigned. The final mask is implicit.
\item
  Per-symbol enumerative coding given a predicted support $S_v$
  containing $M_v$ of size $s_v$: encode $M_v$ by its rank among
  $\binom{s_v}{k_v}$ possibilities, with an escape mode for positions
  in $M_v \setminus S_v$.
\item
  Joint enumerative coding under the multinomial: encode the partition
  by its rank among the $M(N; k_0, \ldots, k_{15})$ valid partitions
  with the given counts.
\end{itemize}

\textbf{5.5 Baseline, lemma, and the corrected gain theorem}

A natural but incorrect baseline for the partition residual is the
per-symbol enumerative sum $\sum_v \log_2 \binom{N}{k_v}$. This
overcounts: the 16 subset choices it sums over are not constrained to
be disjoint, so the same position may be assigned to multiple $v$ in
the enumeration, and the resulting count exceeds the true number of
valid partitions. The correct baseline is the multinomial, which we
state in Theorem 5.2 below; the gain theorem in Theorem 5.3 is
expressed against it.

\textbf{Lemma 5.1.} \emph{Under the factorized regime, per-position
arithmetic coding under $Q$ and per-position Type-I arithmetic coding
under the equivalent sequential conditional $Q$ are bit-equivalent:
both produce code lengths within $\varepsilon_N$ of
$-\sum_i \log_2 Q_{x_i}(i)$.}

\emph{\textbf{Proof.}} Both coders are arithmetic coders driven by the
same per-position probability vector $Q_*(i)$. Arithmetic coding
redundancy is a constant overhead independent of which axis of the
data the coder is sweeping. The two code lengths therefore agree up to
$\varepsilon_N$. \hfill$\blacksquare$

\textbf{Theorem 5.2 (Multinomial baseline for unconstrained partition
coding).} \emph{Let $N$ be a block length, and let
$(k_0, \ldots, k_{15})$ be the symbol counts with $\sum_v k_v = N$.
Without side information, the number of valid 16-way disjoint
partitions of $\{1, \ldots, N\}$ with these counts is
$M(N; k_0, \ldots, k_{15})$. The enumerative code length for the
partition is therefore $\log_2 M(N; k_0, \ldots, k_{15})$ bits, plus
the cost of describing the count vector itself.}

\emph{\textbf{Proof.}} A 16-way partition of $\{1, \ldots, N\}$ with
cardinalities $(k_0, \ldots, k_{15})$ is in bijection with a length-$N$
labeling of the positions in which label $v$ occurs $k_v$ times. The
number of such labelings is the multinomial coefficient $N!$ divided
by the product over $v$ of $k_v!$ by direct combinatorial counting.
Enumerative coding describes the partition by its rank in this class.
\hfill$\blacksquare$

\textbf{Theorem 5.3 (Type-II support-set advantage with escape).}
\emph{Let $\Pi(X) = (M_0, \ldots, M_{15})$ be the true partition with
counts $k_v$ and predicted support sets $S_v$ of size $s_v$.
Encode symbols in fixed order $v = 0, 1, \ldots, 15$. Let
$U_v = \bigcup_{u < v} M_u$ be the positions already claimed,
$A_v = \{1, \ldots, N\} \setminus U_v$ the remaining universe
($|A_v| = N - \sum_{u<v} k_u$), $S'_v = S_v \cap A_v$ the support
restricted to $A_v$ with $s'_v = |S'_v|$, and
$\alpha_v = |M_v \setminus S'_v|$ the escape count.
After transmitting the count vector $(k_v)$ and the escape-count
vector $(\alpha_v)$, encode $M_v \cap S'_v$ by its rank among
$\binom{s'_v}{k_v - \alpha_v}$ subsets of $S'_v$ of that size, and
$M_v \setminus S'_v$ by its rank among
$\binom{|A_v| - s'_v}{\alpha_v}$ subsets of $A_v \setminus S'_v$ of
that size. The total partition code length is at most}

\[
L_{R\Pi} \leq \sum_v \left[\log_2 \binom{s'_v}{k_v - \alpha_v} + \log_2 \binom{|A_v| - s'_v}{\alpha_v}\right] + L(\text{counts}) + L(\alpha) + \varepsilon_N,
\tag{5.1}
\]
\emph{where one may take
$L(\alpha) \leq \sum_v \lceil \log_2 (k_v + 1) \rceil + O(1)$ bits to
describe the escape-count vector explicitly.}

\emph{Comparing against the multinomial baseline
$\log_2 M(N; k_0, \ldots, k_{15})$ of Theorem 5.2, the positional
gain is}

\[
G_{\text{pos}} = \log_2 M(N; k_0, \ldots, k_{15}) - \sum_v \left[\log_2 \binom{s'_v}{k_v - \alpha_v} + \log_2 \binom{|A_v| - s'_v}{\alpha_v}\right].
\tag{5.2}
\]

\emph{Type-II is conditionally advantageous when $G_{\text{pos}}$
exceeds $L(M_{\text{pos}}) + L(V) + L(\Theta) + L(\alpha) +
\varepsilon_N$ plus the cost of describing the count vector and any
per-symbol support headers.}

\emph{\textbf{Proof.}} The decoder reconstructs $M_0, M_1, \ldots$ in
order. At step $v$ the available universe is $A_v$ (positions not
yet claimed by $M_0, \ldots, M_{v-1}$), so the newly decoded set is
automatically disjoint from all previous $M_u$. The transmitted
$\alpha_v$ specifies the two subset sizes; the two enumerative ranks
then identify $M_v \cap S'_v$ within $S'_v$ and $M_v \setminus S'_v$
within $A_v \setminus S'_v$ uniquely. Summing the two log-binomials
across $v$ and adding the count, escape-count, and arithmetic-coder
overheads gives (5.1). The gain identity (5.2) follows by subtracting
from the multinomial baseline. \hfill$\blacksquare$

Theorem 5.3 differs from a natural but flawed formulation in two
respects. First, it uses the multinomial coefficient as the baseline,
which is the correct entropy of an unconstrained partition with
known counts; the looser $\sum_v \log_2 \binom{N}{k_v}$ baseline
overstates the gain by allowing inconsistent subset choices. Second,
it admits the realistic case where $M_v$ is not entirely contained
in $S_v$ by introducing the per-symbol escape count $\alpha_v$,
which carries its own enumerative cost. The bound becomes
particularly clean when $\alpha_v = 0$ for all $v$: the second
log-binomial vanishes and the formula reduces to
$\sum_v \log_2 \binom{s'_v}{k_v}$ (with the chained $s'_v = |S_v \cap A_v|$,
not the raw $s_v$), recovering the idealized oracle bound on the
remaining universe at each step.

\textbf{5.6 Type-II vs Type-I: a separation theorem}

Lemma 5.1 shows that Type-II under a factorized $Q$ is operationally
identical to per-position Type-I arithmetic coding. The natural
question is whether semi-non-factorized Type-II provides a genuine
separation. The next theorem answers this in the affirmative for a
concrete family of joint distributions.

\textbf{Theorem 5.4 (Type-II semi-non-factorized separation,
amortized).} \emph{There exists a family of distributions $D_N$ over
length-$N$ nibble strings, and a positive constant $c > 0$, such that
for every sufficiently large $N$ the optimal factorized per-position
arithmetic coder driven by the true marginals
$Q^{\text{marg}}_v(i) = \Pr_{X \sim D_N}(x_i = v)$ has expected code
length at least $4N$ bits, while a semi-non-factorized Type-II coder
that has access to a single bit per position of joint side information
(membership in a structured set) achieves amortized expected length at
most $(4 - c)N + O(\log N)$ bits per block, when amortized over
$B_N = \Omega(N^2)$ blocks sharing the same generator-matrix model
header. The amortized per-block separation is therefore $\Theta(N)$.}

\emph{\textbf{Proof.}} Fix a rational parameter $c \in (0, 4)$ and
set $d := \lfloor (4 - c) N \rfloor$ (the integer message
dimension; the floor causes at most a $1/N$ additive loss in $c$
that is absorbed into the $O(\log N)$ slack of the theorem
statement and does not affect the $\Theta(N)$ amortized separation).
Let $G$ be a $d \times 4N$ generator matrix over $\mathrm{GF}(2)$
of a linear code $C_N$ satisfying the following two properties:
\begin{itemize}
\item[\textbf{(P1)}]
for every codeword position $j \in \{1, \ldots, 4N\}$, the $j$-th
coordinate of $mG$ (where $m \in \mathrm{GF}(2)^d$ is the message)
is a non-trivial $\mathrm{GF}(2)$-linear function of $m$ ---
equivalently, the $j$-th column of $G$ is non-zero;
\item[\textbf{(P2)}]
for every nibble position $i \in \{1, \ldots, N\}$, the four
columns of $G$ at codeword coordinates $4(i-1)+1, \ldots, 4i$ are
linearly independent in $\mathrm{GF}(2)^d$;
\item[\textbf{(P3)}]
$G$ has full row rank $d$.
\end{itemize}
Such matrices exist with high probability over the uniform
distribution on $d \times 4N$ matrices over $\mathrm{GF}(2)$. For
\textbf{(P1)}, each individual column is the zero vector in
$\mathrm{GF}(2)^d$ with probability $2^{-d}$; a union bound over the
$4N$ columns gives failure probability $\leq 4N \cdot 2^{-d}$. For
\textbf{(P2)}, four fixed columns are linearly dependent iff one of
the $2^4 - 1 = 15$ non-zero $\mathrm{GF}(2)$-linear combinations
yields the zero vector; each such combination of four uniform
independent random vectors in $\mathrm{GF}(2)^d$ is itself uniform
on $\mathrm{GF}(2)^d$ and vanishes with probability $2^{-d}$. Union
over $N$ nibble positions: $\leq 15 N \cdot 2^{-d}$. For
\textbf{(P3)}, a fixed non-zero row-combination
$\lambda \in \mathrm{GF}(2)^d \setminus \{0\}$ produces
$\lambda G \in \mathrm{GF}(2)^{4N}$ which vanishes identically (on
all $4N$ columns) with probability $2^{-4N}$ over the random choice
of $G$; a union bound over the $2^d - 1 < 2^d$ non-zero combinations
gives failure probability $\leq 2^{d - 4N} = 2^{-cN}$.
Combined failure probability:
$\leq 19 N \cdot 2^{-(4-c)N} + 2^{-cN}$, which is $o(1)$ as
$N \to \infty$ for any fixed $c \in (0, 4)$.

Let $S_N \subset \{0, \ldots, 15\}^N$ be the image of all $2^d$
messages under $G$, with the $4N$ output bits parsed as $N$ nibbles.
Property \textbf{(P3)} implies that $G$ has full row rank $d$, so the
encoding map $m \mapsto mG$ is injective and
$|S_N| = 2^d = 2^{\lfloor (4-c)N \rfloor}$. Define $D_N$ as the uniform distribution
on $S_N$.

\textbf{(i) Marginal uniformity.} Fix a nibble position
$i \in \{1, \ldots, N\}$. The four bits encoding the nibble at
position $i$ are values of four $\mathrm{GF}(2)$-linear forms on the
message $m \in \mathrm{GF}(2)^d$, which by \textbf{(P2)} are linearly
independent. The standard fact that $k$ linearly independent
$\mathrm{GF}(2)$-linear forms on a uniformly random vector in
$\mathrm{GF}(2)^d$ induce the uniform distribution on
$\mathrm{GF}(2)^k$ (the joint map is a surjective
$\mathrm{GF}(2)$-linear homomorphism with kernel of size $2^{d - k}$,
so each image has the same preimage size) then gives: each of the
16 bit patterns at position $i$ occurs with probability exactly
$1/16$. Therefore $Q^{\text{marg}}_v(i) = 1/16$ for all $i, v$.

\textbf{(ii) Factorized lower bound.} The cross-entropy of any
$X \sim D_N$ under the factorized model with $Q_v(i) = 1/16$ is
exactly $N \cdot \log_2 16 = 4N$ bits; this is also the optimal
expected code length of an arithmetic coder driven by
$Q^{\text{marg}}$, up to the redundancy $\varepsilon_N$.

\textbf{(iii) Joint upper bound.} A semi-non-factorized coder that
encodes membership in $S_N$ as the unique $d$-bit message index per
block, followed by deterministic expansion via $G$, achieves per-block
content length
$\log_2 |S_N| + \varepsilon_N = d + \varepsilon_N \leq (4 - c)N +
\varepsilon_N$ bits (with the $\leq$ from the floor in
$d = \lfloor (4-c)N \rfloor$, an additive loss of at most $1$ bit
absorbed into the $O(\log N)$ slack), i.e., savings of at least
$4N - (4-c)N = cN$ bits per block \emph{before} model-header
amortization. Describing the generator matrix $G$ itself costs
$L(G) = d \cdot 4N \leq 4(4-c)N^2$ bits, paid once and absorbed
into the $L(M_{\text{pos}})$ term of (3.1). When the same $G$ is
reused across $B_N$ blocks, the per-block model contribution is
$L(G)/B_N$; taking $B_N = \Omega(N^2)$ (e.g., $B_N = N^2$) makes
this $O(1)$ bits per block, absorbed into the $O(\log N)$ overhead.
For fewer blocks the per-block cost is correspondingly higher.

\textbf{(iv) Separation.} The linear-code joint coder of step (iii)
transmits only the $d$-bit message index per block; the decoder
recovers $X$ by computing $m G$ and parsing as nibbles, with no
auxiliary count vector required (the construction of $D_N$ already
fixes $|S_N| = 2^d$ and the marginals are uniform by step (i)).
Subtracting the per-block upper bound (iii) from the per-block lower
bound (ii) thus gives
\[
\Delta(N) \geq 4N - \left[(4 - c)N + \varepsilon_N + 1\right] = cN - \varepsilon_N - 1 = \Theta(N)
\]
for any fixed $c \in (0, 4)$, once the one-time generator-matrix
description $L(G) \leq 4(4-c)N^2$ is amortized across $B_N = \Omega(N^2)$
blocks (the per-block model contribution then absorbs into the
$O(1)$ overhead). The $-1$ from the floor in $d$ is absorbed.
\hfill$\blacksquare$

Theorem 5.4 establishes that semi-non-factorized Type-II strictly
separates from \emph{factorized marginal} per-position coders: coders
that produce one categorical distribution per position using only
each position's marginal $Q^{\text{marg}}_v(i)$, with no joint
conditioning. A genuine autoregressive coder (PixelRNN, PixelCNN,
NNCP) that uses chain-rule conditionals
$Q(x_i \mid x_1, \ldots, x_{i-1})$ can in principle achieve the
joint entropy $H(D_N) = \lfloor (4 - c)N \rfloor$ on the same $D_N$
and is \emph{not} separated by Theorem 5.4; the relevant boundary is
between models with and without joint conditioning, not between RNN
and CNN architectures. The separation is $\Theta(N)$, not
$\Theta(\log N)$: joint structure beyond per-position marginals can
carry a constant fraction of a bit per position even for the
marginal-conditioning case. The linear-code
family used in the proof is a worst case for factorized models
(marginals are exactly uniform) and a best case for joint coders (the
joint distribution lives on a $2^{\lfloor (4-c)N \rfloor}$-element set). Realistic
data classes interpolate between these extremes; the practical regime
in which Type-II is advantageous is one where the data lives on a
small, structurally describable subset of $\{0, \ldots, 15\}^N$ whose
per-position marginals appear approximately uniform. Count constraints
alone yield only a logarithmic separation (the Stirling correction),
so richer joint structure --- geometric constraints, algebraic
parities, sparsity patterns, sortedness --- is required to exhibit the
$\Theta(N)$ regime.

The linear-code construction of Theorem 5.4 is one extreme of a more
general phenomenon. The next result identifies the exact
information-theoretic quantity governing the separation between
factorized and joint coding for any source distribution, and the
following remark connects this quantity to the manifold hypothesis from
statistical learning theory and topological data analysis.

\textbf{Theorem 5.5 (Multi-information separation).} \emph{Let $D_N$
be any distribution over $A_{16}^N$. Define the total correlation
(multi-information, Watanabe 1960; see also Studen\'y and Vejnarov\'a
1998)
$\mathrm{TC}(D_N) = \sum_{i=1}^N H_{\text{marg}}(X_i) - H(D_N)$, where
$H_{\text{marg}}(X_i)$ is the marginal entropy at position $i$ under
$D_N$. Then:
\begin{enumerate}
\item[\textnormal{(F)}] (Factorized rate.) The optimal factorized
per-position arithmetic coder driven by the true marginals
$Q^{\mathrm{marg}}_v(i) = \Pr_{X \sim D_N}(x_i = v)$ has expected
length \emph{at least} $\sum_i H_{\mathrm{marg}}(X_i)$ bits per block
(Shannon's converse for product-distribution coders), and a block
arithmetic coder under those marginals attains
$\sum_i H_{\mathrm{marg}}(X_i) + \varepsilon_N$ in expectation
(achievability), with $\varepsilon_N$ the per-block entropy-coder
redundancy of §2.3.
\item[\textnormal{(J)}] (Joint rate.) Any uniquely decodable lossless
code for $X$ has expected length \emph{at least} $H(D_N)$ bits per
block (Shannon's converse), and a block arithmetic coder driven by
the joint distribution $D_N$ attains
$H(D_N) + \varepsilon_N$ in expectation (achievability).
\item[\textnormal{(S)}] (Separation.) The gap between the optimal
factorized and the optimal joint expected per-block rates equals
$\mathrm{TC}(D_N) + O(\varepsilon_N)$.
\end{enumerate}}

\emph{\textbf{Proof.}} (F) The converse is Shannon's source coding
theorem applied to the product distribution
$\prod_i Q^{\mathrm{marg}}(\cdot \mid i)$; the achievability uses
the same product model with an arithmetic coder (Cover and Thomas
2006, Theorem 5.4.1), incurring at most $\varepsilon_N$ redundancy
per block (§2.3). The expected length under $D_N$ of any code matched
to this product model is the cross-entropy
$\mathbb{E}_{D_N}[-\sum_i \log_2 Q^{\mathrm{marg}}_{X_i}(i)] =
\sum_i H_{\mathrm{marg}}(X_i)$ (since the marginals match).
(J) The converse is Shannon's source coding theorem applied directly
to $D_N$; the achievability is block arithmetic coding under $D_N$,
again with $\varepsilon_N$ redundancy. Neither argument uses AEP /
typical-set machinery (which would require an iid asymptotic setup
that we do not assume).
(S) Subtracting the two attained rates gives
$\sum_i H_{\mathrm{marg}}(X_i) - H(D_N) = \mathrm{TC}(D_N)$ by
definition. Total correlation is non-negative, with equality iff
positions are mutually independent under $D_N$.
\hfill$\blacksquare$

Theorem 5.5 generalizes Theorem 5.4 from the uniform-marginal
linear-code construction to arbitrary source distributions. The
quantity $\mathrm{TC}(D_N)$ is the classical \emph{multi-information}
of Watanabe (1960), and the gap between factorized and joint coding
is a standard consequence of the chain rule for entropy; the novelty
is the placement of this quantity as the operational separator
between Type-I and Type-II under the RNR framework. The $\Theta(N)$
separation observed in Theorem 5.4 corresponds to
$\mathrm{TC} \geq cN - 1$ for the linear-code family (sum of
marginals $\sum_i H_{\text{marg}}(X_i) = 4N$ minus joint entropy
$H(D_N) = \lfloor (4-c)N \rfloor$). For general
$D_N$, the magnitude of $\mathrm{TC}(D_N)$ determines whether Type-II
provides a constant-factor or vanishing improvement over factorized
coding.

\textbf{Remark 5.6 (Manifold-hypothesis intuition, not a theorem).}
The manifold hypothesis from statistical learning theory and
topological data analysis (Donoho 2000; Tenenbaum, de Silva, Langford
2000; Carlsson 2009) is that high-dimensional natural data ---
images, text, audio, structured records --- concentrates near a
low-dimensional structure in the ambient space. It is tempting to
conclude that if $D_N$'s typical realizations live near a structure
of intrinsic dimension $d \ll N$, then $H(D_N) = O(d \log N)$ and
hence $\mathrm{TC}(D_N) = \Theta(N)$ under near-uniform marginals.
\emph{This implication is not a theorem.} A discrete distribution
supported near a $d$-dimensional manifold can have arbitrarily large
Shannon entropy unless additional structure is assumed: a
quantization scale $\Delta$, a bound on the manifold's reach
(local-curvature radius) and on the density (e.g., Lipschitz density
with respect to the manifold's intrinsic measure), and a bound on
the parameter complexity of the manifold itself. Under explicit such
assumptions one can derive entropy bounds of the form
$H(D_N) \leq d \log_2(1/\Delta) + C(d, \mathrm{reach}, \text{density})$;
without them, low intrinsic dimension is not enough.

The cited empirical estimates also do not support a direct
quantitative claim about whole natural images. Carlsson, Ishkhanov,
de Silva and Zomorodian (2008) found low-dimensional local topology
(specifically a Klein-bottle-like structure) in $3 \times 3$
high-contrast image patches; they do not establish
$d \in [10^2, 10^3]$ for full images in ambient dimension $10^6$.
Levina and Bickel (2004) gave a maximum-likelihood intrinsic-dimension
estimator with known biases on noisy real data. We therefore treat
the manifold connection as informal motivation only. Whether any
specific natural-data class $D_N$ satisfies $\mathrm{TC}(D_N) =
\Theta(N)$ is an empirical question (deferred to §11), and
constructing an efficient enumerative encoder for the corresponding
typical set is open.

\textbf{5.7 Worked numerical example}

Consider a block of $N = 1024$ nibbles with counts $k_v = 64$ for
$v = 0, \ldots, 15$ (i.e., uniformly balanced). The multinomial
baseline of Theorem 5.2 evaluates to
\[
\log_2 \!\left[\frac{1024!}{(64!)^{16}}\right] \approx 4033.08 \text{ bits.}
\]
The plain per-byte raw cost is $4 \cdot 1024 = 4096$ bits. A factorized
arithmetic coder under the uniform marginal pays approximately 4096
bits. The multinomial coder saves 75 bits just from the count
constraint alone.

Suppose a positional predictor narrows each symbol's candidate
support to $s_v \leq 256$ with no escapes ($\alpha_v = 0$).
Theorem 5.3's chained residual requires
$s'_v = |S_v \cap A_v|$, and since $|A_v|$ shrinks by $k_{v-1} = 64$
after each symbol is placed, the chain is
$|A_0|, |A_1|, \dots, |A_{15}| = 1024, 960, \dots, 64$. The
predictor's intended support size $s_v = 256$ is achievable as
$s'_v = 256$ only when $|A_v| \geq 256$, i.e., for $v = 0, \dots,
12$ (with $|A_{12}| = 256$); for $v = 13, 14, 15$ the chain caps
the support at $|A_v| = 192, 128, 64$ respectively. Under the
best-case chaining (the predictor picks $S_v \subseteq A_v$ at each
step), the Type-II residual cost is
\[
\sum_{v=0}^{12} \log_2 \binom{256}{64}
+ \log_2 \binom{192}{64}
+ \log_2 \binom{128}{64}
+ \log_2 \binom{64}{64}.
\]
The four binomials evaluate (exactly, to 4 decimal places) to
$203.5669$, $172.2773$, $124.1714$, and $0$ bits respectively,
giving a total of
$13 \cdot 203.5669 + 172.2773 + 124.1714 + 0 \approx 2942.82$ bits.
Compared to the multinomial baseline of $4033.08$ bits, the
apparent gain is $4033.08 - 2942.82 \approx 1090.27$ bits per block.
Adding a model cost of, say, $200$ bytes $= 1600$ bits amortized
over $m = 8$ blocks gives $200$ bits per block, leaving an
effective gain of approximately $890$ bits per block. This is a
saving of roughly $111$ bytes per $512$-byte equivalent block, a
competitive figure on structured positional data.

\textbf{5.8 Support selection and the arithmetic-coding benchmark}

Theorem 5.3 bounds the Type-II partition residual under a candidate
support $S_v$ of fixed size $s_v$. The remaining design question is
how to choose $S_v$ from the predictor's probability field $Q$. In the
semi-non-factorized regime this is a joint support-optimization
problem, not a problem determined by the one-position marginals alone.

\textbf{Remark 5.7 (Support optimization remains open).}
\emph{A previous version asserted that, under a monotonicity condition
on
\[
A_s(\alpha)
= \log_2 \binom{s}{k_v - \alpha}
+ \log_2 \binom{N - s}{\alpha},
\]
the support of size $s$ minimizing the expected Theorem 5.3 residual
is obtained by taking the top-$s$ positions sorted by $Q_v(i)$. This
assertion is false in general. Concrete counterexample at $N = 7$,
$k = 2$, $s = 2$: let $M_v$ be the random 2-subset of $\{0, \dots, 6\}$
distributed over the five atoms
$\{4,5\}, \{3,4\}, \{0,4\}, \{2,4\}, \{1,2\}$ with probabilities
$0.12, 0.04, 0.36, 0.20, 0.28$ respectively (these sum to 1, and
$A_s(\alpha)$ is non-decreasing on the feasible $\alpha \in \{0,1,2\}$).
The top-$s$-by-$Q_v$ support $\{4, 2\}$ has strictly higher expected
$A_s$ residual than the support $\{4, 0\}$, exhibited in
$\texttt{scripts/verify/prop\_5\_7\_counterexample.py}$ in the
companion repository. Even when $A_s(\alpha)$ is non-decreasing, the
expected value of $A_s(|M_v \setminus S_v|)$ depends on the full
joint law of the random set $M_v$, not only on the marginals
$Q_v(i) = \Pr(i \in M_v)$.}

\emph{Thus the actual fixed-size support-selection problem is}
\[
\min_{S_v \subseteq \{1, \ldots, N\},\, |S_v| = s}
\mathbb{E}_{M_v \sim Q}
\left[
A_s\bigl(|M_v \setminus S_v|\bigr)
\right],
\]
\emph{with the remaining-universe correction of Theorem 5.3 applied
when symbols are encoded sequentially. Determining an optimal support
for a general semi-non-factorized predictor $Q$ is left open. Selecting
the largest marginal values $Q_v(i)$ is a natural and common heuristic,
but it is not a theorem of optimality for the general support-set
residual.}

For comparison, in the factorized regime of Section 5.3, per-position
arithmetic coding under $Q$ has expected length equal to the
cross-entropy
\[
\mathbb{E}_D \!\left[-\sum_i \log_2 Q_{X_i}(i)\right],
\]
which equals the conditional entropy only when the factorized
predictor matches the true conditional law. This arithmetic-coding
benchmark is separate from the support-selection problem above.

\textbf{6. Type-III: Tensor Entropy-Field Coding}

\textbf{6.1 General formulation}

Type-III generalizes Type-II in two directions. First, the
representation target $Z = T(X)$ is a multi-dimensional tensor: it may
index over symbol class, position, scale, context class, or token
family. Second, the model predicts not only a first-order field $Q$
over $Z$ but also a second-order field $U$ interpreted as predicted
uncertainty, repair cost, or coding mode suitability. The decoder,
having both $Q$ and $U$ deterministically, can select coding
strategies tile-by-tile and recover the exact $Z$, hence
$X = T^{-1}(Z)$.

Three subtypes specialize this construction. Type-III-A treats $U$ as
a per-tile coding-mode score and uses it to switch between Type-I,
Type-II, and other coders. Type-III-B treats $Z$ as a byte tensor in
which high and low nibbles are jointly modeled. Type-III-C treats $Z$
as a token-interval cover under a hierarchically indexed learned
dictionary.

\textbf{6.2 Type-III-A: entropy-field meta-prediction}

Type-III-A predicts a per-tile, per-mode coding-cost field $U_m(t)$,
where $t$ indexes a tile of the block (a fixed-size window of 64,
128, 256, or 512 positions) and $m \in M$ indexes a coding mode
(Type-I, Type-II, Type-III-B, Type-III-C, raw escape). The
deterministic selector chooses
$\hat{m}(t) = \operatorname*{arg\,min}_m U_m(t)$, and each tile is
encoded with the selected mode. Because $U$ is a function of
decoder-reproducible state, no per-tile mode flag needs to be stored,
except when the encoder departs from the predicted argmin.

Let $L_m(t)$ denote the actual code length when tile $t$ is encoded
with mode $m$, and $U_m(t)$ the predicted value. Define the per-mode
prediction error $\varepsilon_t = \max_m |U_m(t) - L_m(t)|$. The next
theorem bounds the gap between the cost of the model-selected mode
and the cost of the oracle-best mode.

\textbf{Theorem 6.1 (Approximate-oracle mode selection (Type-III-A)).}
\emph{Let $m^\star(t) = \operatorname*{arg\,min}_m L_m(t)$ be the
oracle best mode for tile $t$,
$\hat{m}(t) = \operatorname*{arg\,min}_m U_m(t)$ the model-selected
mode, and $\varepsilon_t = \max_m |U_m(t) - L_m(t)|$ the per-mode
$L^\infty$ prediction error on tile $t$. Then}
\[
L_{\hat{m}(t)}(t) \leq L_{m^\star(t)}(t) + 2\varepsilon_t.
\tag{6.1}
\]

\emph{\textbf{Proof.}} By the prediction error bound,
$L_{\hat{m}}(t) \leq U_{\hat{m}}(t) + \varepsilon_t$. By the
definition of $\hat{m}$, $U_{\hat{m}}(t) \leq U_{m^\star}(t)$. By the
prediction error bound applied to $m^\star$,
$U_{m^\star}(t) \leq L_{m^\star}(t) + \varepsilon_t$. Chaining:
$L_{\hat{m}}(t) \leq L_{m^\star}(t) + 2\varepsilon_t$.
\hfill$\blacksquare$

The $2\varepsilon$ factor is tight. The theorem says: a Type-III-A
coder loses at most twice the per-mode prediction error relative to an
omniscient mode selector. When the predictor is accurate on per-mode
cost (small $\varepsilon_t$), the loss vanishes. This is the formal
justification for predicting prediction cost.

A useful corollary is that Type-III-A composes with any set of
underlying coders without loss: adding a new mode $m'$ to $M$ can only
decrease the oracle cost, and the $2\varepsilon$ bound continues to
hold provided the predictor is calibrated on $m'$. This gives
Type-III-A its role as the top-level scheduler in any RNR composition.

\textbf{Theorem 6.2 (Tightness of the $2\varepsilon$ bound).}
\emph{For every $\varepsilon > 0$ and every $\delta > 0$, there exist
a tile, a mode set $M$ of size 2, true costs $L_m(t)$, and predictions
$U_m(t)$ with $\max_m |U_m(t) - L_m(t)| \leq \varepsilon$ such that
the model-selected mode $\hat{m}$ satisfies
$L_{\hat{m}}(t) \geq L_{m^\star}(t) + 2\varepsilon - \delta$.}

\emph{\textbf{Proof.}} Set $L_1 = \delta/2$ and
$L_2 = 2\varepsilon - \delta/2$, so $m^\star = 1$ and
$L_{m^\star} = \delta/2$. Set $U_1 = \delta/2 + \varepsilon$
(overestimate by $\varepsilon$) and $U_2 = \varepsilon - \delta/2$
(underestimate by $\varepsilon$). Then $|U_1 - L_1| = \varepsilon$ and
$|U_2 - L_2| = \varepsilon$, both within bound. Since
$U_2 = \varepsilon - \delta/2 < \varepsilon + \delta/2 = U_1$, the
argmin selects $\hat{m} = 2$. Therefore
$L_{\hat{m}} - L_{m^\star} = (2\varepsilon - \delta/2) - \delta/2 =
2\varepsilon - \delta$. As $\delta \to 0$, the gap approaches
$2\varepsilon$. \hfill$\blacksquare$

Theorem 6.2 shows that the $2\varepsilon$ factor in Theorem 6.1
cannot be improved without additional assumptions on the predictor's
error structure (for example, one-sided errors, calibrated quantile
bounds, or bounded total variation between $U$ and $L$). Standard
concentration arguments give better constants in the typical case but
the worst-case factor is exactly 2.

\textbf{6.3 Type-III-B: byte tensor coding}

Type-III-B extends Type-II from nibble to byte alphabets by jointly
modeling the high and low nibbles of each byte. Write byte
$b = 16h + l$ with $h, l \in A_{16}$. The model predicts a
$16 \times 16 \times N$ tensor $Q_{h,l}(i)$ approximating
$\Pr(h_i = h, l_i = l \mid C)$, normalized so $\sum_{h,l} Q_{h,l}(i) =
1$ for every $i$. The construction is exactly Type-II with a 256-symbol
alphabet but with the alphabet structure exposed.

Two coding strategies are available. Direct partition coding treats
the alphabet as 256 flat classes and applies Theorem 5.3 with
$A_{256}$. Hierarchical coding decomposes the byte partition along the
high-nibble axis: first encode the high-nibble partition
$M^H_h = \{ i : h_i = h \}$, then for each $h$ encode the low-nibble
conditional partition
$M^{L \mid H = h}_l = \{ i \in M^H_h : l_i = l \}$. The hierarchical
scheme is preferred when the predictor decomposes naturally as
$Q_{h,l}(i) = Q^H_h(i) \cdot Q^{L \mid H}_l(i, h)$.

\textbf{Theorem 6.3 (Type-III-B mutual-information advantage).}
\emph{Let $H$, $L$, and $C$ denote the random variables corresponding
to the high nibble, low nibble, and public context at a fixed
position. Under an ideal entropy coder, independent nibble coding pays
expected length $H(H \mid C) + H(L \mid C)$ per position, while joint
byte coding pays $H(H, L \mid C)$. The advantage of byte coding over
independent nibble coding is}

\[
\Delta_B = H(H \mid C) + H(L \mid C) - H(H, L \mid C) = I(H ; L \mid C),
\tag{6.2}
\]

\emph{the conditional mutual information between the high and low
nibble given context, in bits per position. For a block of $N$
positions, Type-III-B is conditionally advantageous over Type-II at
the nibble alphabet whenever the total saving
$\sum_{i=1}^{N} I(H_i ; L_i \mid C_i)$ exceeds the additional model
and metadata cost
$L(M_{\mathrm{III\text{-}B}}) - L(M_{\mathrm{II}}) + L(\Theta_{\mathrm{III\text{-}B}}) - L(\Theta_{\mathrm{II}})$
of the byte-tensor predictor (both sides in bits, both summed over
the block); equivalently, when amortized per position, whenever
$\overline{I(H;L \mid C)}$ exceeds
$[L(M_{\mathrm{III\text{-}B}}) - L(M_{\mathrm{II}}) + L(\Theta_{\mathrm{III\text{-}B}}) - L(\Theta_{\mathrm{II}})] / N$
in bits per position.}

\emph{\textbf{Proof.}} By the chain rule of entropy,
$H(H, L \mid C) = H(H \mid C) + H(L \mid H, C)$. Subtracting from
$H(H \mid C) + H(L \mid C)$ gives $H(L \mid C) - H(L \mid H, C)$,
which by definition is $I(H ; L \mid C)$. Both terms are non-negative,
so $\Delta_B \geq 0$, with equality iff $L$ is conditionally
independent of $H$ given $C$. \hfill$\blacksquare$

The mutual information $I(H ; L \mid C)$ is large for many practical
data classes. ASCII text concentrates on high nibbles 2 through 7 with
structured low-nibble distributions per high nibble; opcode streams
restrict low nibbles depending on the high nibble; pixel data has
locally correlated joint byte distributions. On such data, Type-III-B
can recover several tenths of a bit per byte over Type-II at the
nibble alphabet.

\textbf{Theorem 6.4 (Hierarchical decomposition is
information-theoretically invariant).} \emph{Let $T$ be any rooted
binary tree whose leaves are the 256 byte values, defining a sequence
of nested binary partitions of the alphabet. Encoding a byte $b$ under
$T$ by traversing root-to-leaf and encoding the binary decision at
each node with the \emph{exact induced conditional probability} of
that decision under $P(B \mid C)$ given the prior decisions on the
path, with an ideal sequential entropy coder (e.g.\ arithmetic
coding in the asymptotic-redundancy limit), achieves expected
\emph{ideal} length $H(B \mid C)$ per byte. The expected ideal cost
is independent of the tree shape $T$; it depends only on the joint
distribution $P(B \mid C)$. Practical prefix-code constructions
(Huffman, etc.) achieve $H(B \mid C) + O(1)$ per byte; arithmetic
or ANS coders achieve $H(B \mid C) + o(1)$ amortized over many
bytes (Witten--Neal--Cleary 1987; Duda 2009). \emph{Caveat:} this
invariance fails if the decisions are coded with fixed branch
probabilities (e.g., 1/2 each) instead of the induced conditional
--- in that case the depth of $T$ dominates, and a left-skewed tree
pays substantially more than a balanced one. The theorem applies
only to coders that use the \emph{exact} induced conditional
probability at each node.}

\emph{\textbf{Proof.}} Fix any rooted binary tree $T$ on the 256
byte values. For each byte value $b$, let $\pi_T(b) = (z_1, z_2,
\ldots, z_{\ell_T(b)}) \in \{0,1\}^\star$ denote the unique
root-to-leaf path on $T$, of $T$-dependent depth $\ell_T(b)$ (a
path-length integer, not a code-length). The map $b \mapsto
\pi_T(b)$ is a bijection onto a prefix-free set of binary strings
(the leaves of $T$), so
$\Pr(\pi_T(B) = (z_1, \dots, z_{\ell_T(B)}) \mid C) = \Pr(B = b \mid C)$
for the unique $b$ whose path is $(z_1, \dots, z_{\ell_T(B)})$.
Encoding $b$ by sequentially passing the bits of $\pi_T(b)$ to an
ideal sequential entropy coder (Shannon-Fano arithmetic, range, or
ANS in the asymptotic-redundancy limit), with each bit's coding
probability set to its \emph{exact induced conditional}
$\Pr(z_l = 1 \mid z_1, \ldots, z_{l-1}, C)$, gives an
ideal-coder code length of
$-\log_2 \Pr(\pi_T(B) \mid C) = -\log_2 \Pr(B \mid C)$ bits per
byte. Taking expectation under $B \sim P(\cdot \mid C)$:
\[
\mathbb{E}_B[\text{ideal code length}] = \mathbb{E}_B[-\log_2 \Pr(B \mid C)] = H(B \mid C),
\]
which is independent of $T$. (Note that the \emph{path length}
$\mathbb{E}_B[\ell_T(B)]$ does depend on $T$, and is in general
much larger than $H(B \mid C)$ for unbalanced trees; the
invariance is the expected entropy-coded length, not the expected
path length.) \hfill$\blacksquare$

The practical implication is that the choice of decomposition tree
for Type-III-B (high-low, or any other binary split) affects
encoder/decoder complexity and predictor calibration, but not the
information-theoretic code length. The optimal tree for a specific
practical entropy coder is the one whose binary decisions are best
predicted by the available $Q$; this is a calibration question, not a
coding-rate question. In particular, the canonical high-low
decomposition is optimal when the predictor factorizes naturally as
$Q_{h,l}(i) = Q^H_h(i) \cdot Q^{L \mid H}_l(i, h)$, which is typical
for transformer predictors trained on byte streams.

\textbf{6.4 Worked numerical example for Type-III-B}

Consider an ASCII-text block where the empirical marginal entropies
are $H(H \mid C) = 2.4$ bits and $H(L \mid C) = 2.9$ bits per byte,
while the joint entropy is $H(H, L \mid C) = 4.7$ bits. Independent
nibble coding pays $2.4 + 2.9 = 5.3$ bits per byte; joint byte coding
pays $4.7$ bits per byte. The mutual-information advantage is
$\Delta_B = 5.3 - 4.7 = 0.6$ bits per byte, or 0.075 bytes per byte
(about 7.5\% saving) before model and metadata overhead. Over a
4096-byte block, this is 307 bytes of raw saving against independent
nibble coding, against an extra model cost that is typically a few
tens of bytes per amortized block.

\textbf{6.5 Type-III-C: token tensor coding}

Type-III-C generalizes Type-III-B from the byte alphabet to a learned
token dictionary $D = \{\tau_0, \ldots, \tau_{K-1}\}$. Tokens may be
single bytes, byte pairs, longer byte fragments, structural markers,
or --- most generally --- symbolic compositions in which a token's
definition references other tokens in the dictionary. The latter case
makes the dictionary a \emph{straight-line grammar}: a context-free
grammar with one production per non-terminal and a directed acyclic
dependency graph (Kieffer and Yang 2000; Charikar et al. 2005).
Practical dictionaries used by BPE, WordPiece, and MDL-greedy fall
into the flat case (no recursion); the framework also admits the
recursive case, and we use it in the complexity arguments of Theorems
6.7 and 6.8 below.

The model predicts a position-and-scale tensor $Q(k, i, s)$
approximating $\Pr(\tau_k \text{ begins at position } i \text{ at
scale } s \mid C)$, where $s$ is a length or scale class.

The dictionary is structured hierarchically as a grid of $K = a \cdot
b$ classes, with $\text{token\_id} = a \cdot \text{family} +
\text{variant}$. For $K = 1024$, the canonical grid is
$32 \times 32$; for $K = 4096$, the grid is $64 \times 64$. The
hierarchical factorization
$Q(k, i, s) = P(\text{family} \mid i, C) \cdot
P(\text{variant} \mid \text{family}, i, C)$
reduces the output dimension of the predictor at each position from
$K$ to $a + b$, which is asymptotically a quadratic reduction in
softmax cost.

The $32 \times 32$ grid does not span all 65536 byte pairs. It is an
MDL-selected dictionary of 1024 useful fragments. Admission of a token
to $D$ is governed by the standard MDL token gain criterion,
formalized in the next theorem.

\textbf{Theorem 6.5 (Type-III-C marginal-MDL token admission, all
else frozen).} \emph{Let $\tau$ be a candidate token of byte length
$\ell$. Fix a parser that produces a non-overlapping token cover of
the corpus and let $f$ be the \emph{realized} count of $\tau$ in
that parse, restricted to positions currently parsed as raw
literals (i.e., not already covered by any token already in $D$).
This count can be strictly less than the raw occurrence count when
$\tau$ self-overlaps; e.g., for $\tau = $ \texttt{aa} in the
string \texttt{aaa}, the raw count is $2$ but $f = 1$. Positions
where $\tau$ would overlap an existing token are excluded from $f$
by the frozen-parser assumption below. Let
$\mathrm{code\_cost}(\tau)$ be the average bit length to encode an
occurrence of $\tau$ under the predictor and entropy coder,
$\mathrm{dictionary\_cost}(\tau)$ the bit length to describe $\tau$
in the dictionary header, and $\mathrm{parse\_overhead}(\tau)$ the
additional per-occurrence boundary or interval bookkeeping
introduced by including $\tau$. \emph{Crucially, assume that the
parser, the entropy-coder state, all other tokens' code costs,
dictionary costs, parse overheads, and the realized counts of all
other tokens are FROZEN during the admission test for $\tau$
(``ceteris paribus''); the rule evaluates only the marginal effect
of adding $\tau$ to an otherwise fixed dictionary. Adding $\tau$ in
practice may change the parse and the marginal costs of competing
tokens; the test below ignores those interactions, so it gives a
greedy admission decision rather than a globally optimal one.} The
net change in total Type-III-C code length under these frozen-rest
assumptions is}
\[
\Delta L(\tau) = \mathrm{dictionary\_cost}(\tau) + f \cdot \left[\mathrm{code\_cost}(\tau) + \mathrm{parse\_overhead}(\tau) - 8\ell\right].
\]
\emph{Consequently, the per-token greedy MDL admission policy admits
$\tau$ if and only if $\Delta L(\tau) < 0$, equivalently}
\[
f \cdot \left[8\ell - \mathrm{code\_cost}(\tau)\right] > \mathrm{dictionary\_cost}(\tau) + f \cdot \mathrm{parse\_overhead}(\tau).
\tag{6.3}
\]

\emph{\textbf{Proof.}} Without $\tau$ in $D$, the $f$ occurrences of
the byte pattern are coded as raw literals at $f \cdot 8\ell$ bits
total, with no dictionary contribution. With $\tau$ in $D$, the $f$
occurrences cost $f \cdot \mathrm{code\_cost}(\tau)$ for encoding plus
$f \cdot \mathrm{parse\_overhead}(\tau)$ for bookkeeping, and the
dictionary pays a one-time $\mathrm{dictionary\_cost}(\tau)$. The
difference of the two totals is $\Delta L(\tau)$. By the MDL
principle (Rissanen 1978), the admission policy that minimizes total
description length includes $\tau$ iff this difference is negative.
Rearranging gives (6.3). \hfill$\blacksquare$

The token MDL criterion is the standard greedy dictionary-learning
admission rule. Its placement here is to make explicit that Type-III-C
does not store every possible token, only those whose per-corpus
savings cover their description. The criterion is greedy in that it
evaluates each candidate $\tau$ marginally, assuming the rest of the
dictionary is fixed; it is not a global MDL optimum, since admitting
$\tau$ can change the optimal parse and the
$\mathrm{code\_cost}$ of other tokens. The closest computational
counterpart we prove is the restricted normal-form dictionary-only
problem of Theorems 6.7 and 6.8 below, which is shown to be NP-hard
and APX-hard with constant $\gamma_{\mathrm{SGP}} > 1$; the fully
global Type-III-C problem (with model, dictionary, parser, and
repair stream all free) is conjectured at least as hard but its
formal hardness is open (§13.2). For a corpus encoded across many
archives,
$\mathrm{dictionary\_cost}$ is amortized and the greedy admission
criterion grows more permissive.

When tokens are continuous latent variables of a variational model
rather than discrete byte fragments, the appropriate encoding tool is
bits-back coding (Townsend, Bird, Barber 2019). The encoder uses an
auxiliary bit string to sample the latent, encodes the latent and the
data conditional on the latent via ANS, and the decoder recovers both
and extracts the auxiliary bits. The effective rate is the marginal
$H(X)$, not the joint $H(X, Z)$. The next theorem makes this rigorous
and shows how bits-back fits into the $L(R \mid Y)$ decomposition of
§3.3.

\textbf{Theorem 6.6 (Bits-back integration into the RNR
decomposition).} \emph{Let $z$ be a latent variable of a variational
model with prior $P(z \mid Y)$, conditional $P(X \mid z, Y)$, and
variational posterior $q(z \mid X, Y)$. \emph{Quantization
convention for continuous latents:} when $z$ is continuous, fix a
finite-precision discretization $\Delta$ of the latent space common
to encoder and decoder (e.g., $b$-bit fixed-point per coordinate as
used by Townsend-Bird-Barber's BB-ANS implementation), and read all
densities below as probability masses on the discretization cells.
With this convention $-\log_2 P(z \mid Y)$ and $-\log_2 q(z \mid X,
Y)$ are finite code lengths and the dictionary case is the trivial
discretization where $z$ takes values in a discrete set of token
identifiers (so $\Delta$ disappears). Under
bits-back ANS coding with auxiliary bits used to sample $z$ from $q$,
the effective Type-III-C repair length is
$L(R \mid Y) = L_{\text{prior}}(z \mid Y) + L_{\text{cond}}(X \mid z, Y) - L_{\text{aux}}(z)$,
where the three terms encode the latent under the prior, the data
under the conditional, and recover the auxiliary bits respectively.
Its expectation under $X \sim D$ and $z \sim q$ satisfies}
\[
\mathbb{E}[L(R \mid Y)] = \mathbb{E}_{X \sim D}[-\log_2 P_X(X \mid Y)]
+ \mathbb{E}_{X \sim D} \mathrm{KL}\!\left(q(z \mid X, Y) \,\Vert\, P(z \mid X, Y)\right),
\tag{6.4}
\]
\emph{where $P_X(\cdot \mid Y)$ is the model marginal over $X$ given
$Y$. Equivalently, the first term equals
$H_D(X \mid Y) + \mathrm{KL}(D(\cdot \mid Y) \Vert P_X(\cdot \mid Y))$,
the cross-entropy of the data under the model. The bits-back rate
equals the conditional entropy $H_D(X \mid Y)$ only when the model
marginal matches the data conditional ($D = P_X$) and $q$ equals the
true posterior; under both conditions matching the Slepian-Wolf lower
bound. Type-III-C with quantized continuous latents is therefore a strict
generalization of Type-III-C with a discrete dictionary, with the
dictionary case recovered when the discretization $\Delta$ identifies
$q$ with a point mass on a single discrete token.}

\emph{\textbf{Proof.}} Under the quantization convention $\Delta$,
the three-stream bits-back coder transmits
$\log_2 \tfrac{1}{P(z \mid Y)}$ bits for the latent under the
discretized prior, $\log_2 \tfrac{1}{P(X \mid z, Y)}$ bits for the
data under the conditional, and recovers
$\log_2 \tfrac{1}{q(z \mid X, Y)}$ bits from the auxiliary stream
(used to sample $z$). Net code length:
\[
\begin{aligned}
&-\log_2 P(z \mid Y) - \log_2 P(X \mid z, Y) + \log_2 q(z \mid X, Y) \\
&\quad= -\log_2 P(X, z \mid Y) + \log_2 q(z \mid X, Y) \\
&\quad= -\log_2 P(X \mid Y) + \log_2 \!\left[\frac{q(z \mid X, Y)}{P(z \mid X, Y)}\right],
\end{aligned}
\]
using the joint factorization
$P(X, z \mid Y) = P(X \mid Y) \cdot P(z \mid X, Y)$. Taking
expectation under $q(z \mid X, Y)$ conditional on $X$:
$\mathbb{E}[\text{net} \mid X] = -\log_2 P(X \mid Y) + \mathrm{KL}(q(z \mid X, Y) \,\Vert\, P(z \mid X, Y))$.
Taking expectation under $X \sim D$ gives (6.4). The first term is
cross-entropy under the model marginal $P_X(\cdot \mid Y)$, not
entropy unless $P_X(\cdot \mid Y) = D(\cdot \mid Y)$; the second
term is the variational gap. When both vanish the rate is
$H_D(X \mid Y)$. \hfill$\blacksquare$

\textbf{Theorem 6.7 (Restricted Type-III-C dictionary learning is
NP-hard).} \emph{Consider the restricted Type-III-C problem in which
$M$ is fixed to the trivial constant predictor, the repair stream is
the distinguished start token, and $D$ is a recursive straight-line
dictionary in normal form (defined below), with objective the
symbol-count $L_{\mathrm{sym}}(D) = \sum_i |\alpha_i|$. Finding the
$D$ minimizing $L_{\mathrm{sym}}$ subject to $\tau_0$ expanding to
the corpus $X$ is NP-hard.}

\emph{The unrestricted joint problem ($M$, $D$, repair stream all
free, total code length $L(M) + L(D) + L(R\mid Y)$) is conjectured
to be NP-hard but is not established by the present reduction:
dropping the restrictions changes both the feasible set and the
objective function, and for minimization with a larger feasible set
the optimum can be strictly smaller, so a polynomial-time
unrestricted solver does not automatically determine the restricted
optimum. Designing a reduction that forces the unrestricted optimum
to coincide with the restricted optimum is left open.}

\emph{\textbf{Proof.}} Reduce from the Smallest Grammar Problem (SGP)
of Charikar, Lehman, Liu, Panigrahy, Prabhakaran, Sahai and Shelat
(2005). An instance of SGP is a string $s \in \Sigma^\star$ (we may
take $\Sigma = \{0,1\}^8$ for concreteness); the task is to find a
straight-line grammar $G$ with $L(G) = \{s\}$ minimizing the
\emph{grammar size}
$|G| = \sum_{A \to \alpha \in P} |\alpha|$. SGP is NP-hard (Charikar
et al. 2005, Theorem 3.4).

Use \emph{normal-form} dictionaries: $\tau_0$ is the distinguished
start token; every token is reachable from $\tau_0$; the dependency
graph is acyclic; no token has a unit-production right-hand side; and
duplicate terminal expansions are identified. These are the standard
harmless normalizations for straight-line grammars: unreachable
rules, unit productions, and duplicate nonterminals can be removed
without increasing grammar size.

Given an SGP instance $s$, set $X := s$, fix $M$ to the trivial
predictor, and require $\tau_0$ to expand deterministically to $s$.
Define the hardness objective
\[
L_{\mathrm{sym}}(D) = \sum_i |\alpha_i|,
\]
where $\alpha_i$ is the right-hand side of token $\tau_i$ and each
terminal or token reference counts as one symbol. Normal-form
recursive dictionaries on $\Phi(s)$ are in bijection with normal-form
straight-line grammars deriving exactly $\{s\}$ --- map each token to
the homonymous non-terminal and vice versa --- and the bijection
preserves the objective value exactly:
\[
L_{\mathrm{sym}}(D(G)) = |G|.
\]
Thus an exact polynomial-time optimizer for this restricted
Type-III-C case would exactly solve SGP. Since SGP is NP-hard
(Lehman and Shelat 2002; Charikar et al.\ 2005), the restricted
Type-III-C problem is NP-hard. The reduction is computable in linear
time in $|s|$.

The unrestricted-case claim is intentionally not proved by inclusion:
the restricted feasible region is a strict subset of the unrestricted
region, but the objective function also changes (the restricted case
uses $L_{\mathrm{sym}}$ while the unrestricted case uses
$L(M) + L(D) + L(R\mid Y)$), so the standard ``larger feasible set
$\Rightarrow$ smaller optimum'' minimization fact does not let an
unrestricted solver recover the restricted optimum. We leave the
formal unrestricted hardness as an open problem.
\hfill$\blacksquare$

Theorems 6.6 and 6.7 together delineate one slice of the
achievability boundary in Type-III-C. Bits-back gives the optimal
rate for any joint (model, posterior) pair (Theorem 6.6); finding
the optimal pair in the joint (model, dictionary, repair) problem
is not directly proven intractable by Theorem 6.7, which establishes
NP-hardness only for the restricted normal-form dictionary-only
case under the symbol-count objective $L_{\mathrm{sym}}$.
Conjecturally the unrestricted joint problem is at least as hard,
but a formal hardness reduction is open (see §13.2). Practical
implementations nonetheless use heuristic dictionary learning (BPE,
MDL-greedy, neural codebook learning) and amortized variational
inference, both of which sacrifice optimality for tractability.

\textbf{Theorem 6.8 (APX-hardness of joint
dictionary learning, restricted case).} \emph{The restricted
Type-III-C problem of Theorem 6.7 (constant predictor, recursive
normal-form dictionary, symbol-count objective $L_{\mathrm{sym}}$)
is APX-hard, conditional on the established APX-hardness of the
Smallest Grammar Problem: there exists a constant
$\gamma_{\mathrm{SGP}} > 1$ such that no polynomial-time algorithm
achieves approximation ratio better than $\gamma_{\mathrm{SGP}}$ on
SGP unless $P = NP$ (Lehman and Shelat 2002; Charikar et al.\ 2005),
and the bijection of Theorem 6.7 transfers this hardness to the
restricted Type-III-C objective with the same constant. The
APX-hardness of the unrestricted joint problem is left open for the
same reason as in Theorem 6.7: the feasible set is strictly larger
and the objective function differs, so neither inclusion nor a
direct gap-preserving reduction is available.}

\emph{\textbf{Proof.}} We show that the reduction
$\Phi: s \mapsto \Phi(s)$ constructed in Theorem 6.7 preserves
approximability up to constant factors, then transfer the
APX-hardness of SGP.

\emph{Symbol-count normal form.} For the APX statement we use the
normal-form symbol-count objective of Theorem 6.7,
$L_{\mathrm{sym}}(D) = \sum_i |\alpha_i|$. The dictionary-to-grammar
map and grammar-to-dictionary map preserve this objective exactly, so
an $\alpha$-approximation for the restricted Type-III-C problem yields
an $\alpha$-approximation for SGP with no additive bookkeeping slack.

\emph{Transfer of APX-hardness.} Lehman and Shelat (2002) and the
expanded line of work culminating in Charikar, Lehman, Liu,
Panigrahy, Prabhakaran, Sahai and Shelat (2005) establish that SGP
is APX-hard: some explicit constant $\gamma_{\mathrm{SGP}} > 1$
exists such that no polynomial-time algorithm achieves approximation
ratio smaller than $\gamma_{\mathrm{SGP}}$ unless $P = NP$. The
proof routes through gap-preserving reductions from standard
NP-hard problems (the exact source-problem and explicit numerical
value of $\gamma_{\mathrm{SGP}}$ differ between the conference and
journal versions of the cited papers; we do not pin down the
specific numeric constant here). Combined with the exact bijection
above, a polynomial-time $\alpha$-approximation for the restricted
Type-III-C problem with $\alpha < \gamma_{\mathrm{SGP}}$ would yield
a polynomial-time approximation to SGP below the established
threshold, contradicting $P \neq NP$. The unrestricted joint
problem is not addressed by this argument: as in Theorem 6.7, the
restricted feasible region is strict but the objective also
changes, so an unrestricted approximation does not project back to
a restricted-objective approximation. We leave the unrestricted
APX-hardness open. \hfill$\blacksquare$

\textbf{Remark on the constant.} The numerical value of
$\gamma_{\mathrm{SGP}}$ in the SGP literature is small (close to 1)
and depends on the specific reduction used; we do not assert a
specific figure here. The constant $\gamma_{\mathrm{SGP}}$ inherited
by Theorem 6.8 applies to the \emph{restricted} symbol-count
objective only; we do not claim that any constant $\gamma > 1$
inherits to the unrestricted bit-cost objective, since the latter
problem's APX-hardness is itself open (see §13.2). Pinning down a
specific numerical value of $\gamma_{\mathrm{SGP}}$ for any
particular encoding convention, and verifying which source paper
establishes which exact constant, is left to follow-up work.

Together, Theorems 6.7 and 6.8 confirm that the heuristic nature of
practical dictionary learners (BPE, WordPiece, MDL-greedy variants)
is not a deficiency of those algorithms but a necessary consequence of
complexity-theoretic limits. Even with unbounded compute applied to
the dictionary search, polynomial-time algorithms cannot close the
approximation gap unless $P = NP$.

\textbf{7. Composition and mode selection}

RNR types compose by hierarchical mode selection. The top-level
scheduler is a Type-III-A coder whose mode set $M$ contains the
available lower-level coders, including all Type-I, Type-II, and
Type-III-B/C constructions. For each tile, the scheduler predicts a
per-mode cost vector $U_*(t)$, selects the argmin mode, and delegates
encoding to that mode.

\textbf{Theorem 7.1 (Composition correctness, with explicit framing).}
\emph{Let $M$ be a finite set of RNR mode coders, each producing a
uniquely decodable repair stream conditional on its mode flag.
Assume further that the per-tile repair streams are \emph{framed} ---
either each carries an explicit byte-length prefix in the metadata
$L(\Theta)$, or the per-tile encoding is itself prefix-free as a
binary string --- so that the decoder, after parsing the mode flag,
can isolate the bits belonging to tile $t$'s repair stream before
applying the mode-specific decoder. Under this framing assumption,
the Type-III-A composition with mode set $M$, deterministic mode
selection $\hat{m}(t) = \operatorname*{arg\,min}_m U_m(t)$ (with
fixed lexicographic tie-breaking from Definition 3.3), and a
uniquely decodable mode-flag stream produces a uniquely decodable
codeword for $X$.}

\emph{\textbf{Proof.}} The decoder reads the metadata (including the
per-tile length prefixes, if present, that fix the bit boundaries
between successive tiles), reconstructs $U_*(t)$ deterministically
from public state, recovers $\hat{m}(t)$ by the same argmin, isolates
the repair-stream bits belonging to tile $t$ using the framing
information, and applies the decoder of mode $\hat{m}(t)$ to recover
the tile. Each step is uniquely determined; the concatenation is
therefore an injection from $X$ to codewords. \emph{Caveat:} without
the framing assumption, per-mode unique decodability alone is
insufficient: e.g., mode-A codewords $\{0, 01\}$ and mode-B
codewords $\{10, 0\}$ are each uniquely decodable, but the
concatenation $010$ over a mode-A then mode-B tile parses as
$0 \mid 10$ or $01 \mid 0$, ambiguously. The framing assumption
removes this ambiguity by fixing tile boundaries explicitly.
\hfill$\blacksquare$

Composition implies an approximate-oracle bound (not a strict
no-regret guarantee): by Theorem 6.1, including every coder in $M$
makes the Type-III-A composition's total cost at worst within
$2 \sum_t \varepsilon_t$ of the per-tile oracle (which picks the
true argmin mode per tile). This is also at worst within
$2 \sum_t \varepsilon_t$ of any fixed single coder $m^\star$ on the
same data, since the per-tile oracle dominates any fixed single
coder pointwise. When the per-tile best mode varies non-trivially across tiles
(i.e., no single coder is among the per-tile argmin set on every
tile), the per-tile oracle is strictly better than any fixed single
coder; under ties (some tile has multiple argmin modes including
$m^\star$) the oracle and $m^\star$ may agree on cost. The
composition inherits a strict improvement only insofar as
$2 \sum_t \varepsilon_t$ remains below that gap, which is itself an
empirical property of the prediction calibration and is not
guaranteed in general.

\textbf{7.2 Runtime parameter selection}

Several parameters in the framework are input-dependent: the codeword
rate $k$ in Type-I (Code-$k$), the alphabet size $K$ in Type-II, the
support threshold $\tau_v$ for the Type-II support sets (see Remark 5.7
for the open-problem status of optimizing it), the tile size in
Type-III-A, the dictionary size in Type-III-C, and others. No single
setting is optimal across all inputs; the right approach is input
probing at encode time, where the encoder scans the block, estimates
the relevant structural parameter, and selects the parameter
accordingly. The parameter choice is recorded in the metadata stream
$L(\Theta)$ so the decoder can reproduce the selection.

Parameter selection is structurally identical to mode selection: each
parameter value defines a distinct coder, and the Type-III-A scheduler
chooses among them by predicted cost. The composition theorem
(Theorem 7.1) and the approximate-oracle bound (Theorem 6.1) both
extend to parameterized modes without change, with the per-mode
prediction error $\varepsilon_t$ now defined over the cross product of
(mode, parameter) pairs. The cost is the metadata overhead of encoding
the parameter choice, typically $O(\log K_{\text{params}})$ bits per
tile or per block, which is negligible against the per-mode repair
stream.

Two practical patterns are standard. First, batch probing: scan a
small prefix of the block (e.g., the first 1 KiB), estimate key
statistics (entropy rate, marginal sparsity, expected Hamming weight
of Code-$k$ errors), select parameters once, and encode the rest of
the block with the chosen parameters. Second, tiled probing: at each
tile boundary in Type-III-A, re-estimate parameters from the most
recently verified subblocks and adjust. Both patterns fit the
existing framework: the parameter-selection function is itself part
of $M$ (deterministic, decoder-reproducible from public state), and
the parameter values are part of $L(\Theta)$.

This framing resolves the input-dependent design questions that would
otherwise appear as open problems: there is no single optimal
Code-$k$ or optimal alphabet size, only an optimal
parameter-selection function for a given data distribution, and the
latter is part of $M$ just like any other learned component.

\textbf{8. No universal dominance}

\textbf{Theorem 8.1 (No lossless universal compression).} \emph{No
injective map $C$ from $\{0,1\}^n$ to binary strings of length less
than $n$ assigns every input string of length $n$ a strictly shorter
codeword.}

\emph{\textbf{Proof.}} There are $2^n$ binary strings of length $n$.
There are $2^n - 1$ binary strings of length strictly less than $n$,
summing $\sum_{k=0}^{n-1} 2^k$. Any injection from $2^n$ inputs to
fewer than $2^n$ outputs is impossible by the pigeonhole principle.
\hfill$\blacksquare$

Theorem 8.1 is standard and short; we include it to fix the
framework's honest claim. RNR does not assert that any fixed coder
dominates any baseline on every input. The claim is conditional and
distributional: for data classes whose conditional entropy under
deterministic side information is below the baseline rate by more than
the model and overhead cost, RNR is advantageous in expectation.

\textbf{9. Verification and error containment}

Autoregressive RNR modes update the predictor context with previously
decoded subblocks. An undetected decoding error in subblock $i$ would
propagate into the context for subblock $i + 1$ and beyond, with
cascading consequences. We discipline this by requiring that the
decoder update its context only after verifying the candidate decoded
subblock against a stored hash.

Let $X = (X_1, \ldots, X_m)$ be partitioned into verified subblocks,
with stored verification values $v_i = H_v(X_i)$. The decoder produces
$\hat{X}_i$, checks $H_v(\hat{X}_i) = v_i$, and only on success uses
$\hat{X}_i$ to update $C$ for subblock $i + 1$. The following theorem
rules out cross-subblock error propagation.

\textbf{Theorem 9.1 (No cross-subblock context propagation, with
explicit framing).} \emph{Assume the archive header stores, for each
subblock $i$, an explicit byte/bit-offset into the repair stream
(framing), so that subblock $i + 1$'s repair-stream slice is parsed
from a fixed position independent of the bytes recovered in
subblock $i$. Suppose the decoder uses the following
\emph{verified-only context-update rule}: after producing the
candidate $\hat{X}_i$, the decoder evaluates
$H_v(\hat{X}_i) \stackrel{?}{=} v_i$, and only on equality does it
extend the predictor context $C$ with $\hat{X}_i$ for predicting
subblock $i + 1$; on mismatch the decoder aborts (or marks the
remainder unverified and does not update $C$). The predictor input
for subblock $i+1$ is therefore a function only of (a)~the
verified-bytes context $C$ built from $\{\hat{X}_j : j \leq i,
H_v(\hat{X}_j) = v_j\}$ together with any public boundary state in
the header; the rule permits but does not require a hard
\emph{context-reset boundary} (i.e., dropping the verified-bytes
context at each subblock break). Finally, suppose $H_v$ is a
cryptographic hash with negligible collision probability $\delta$
on the considered set of candidate subblocks. Then the probability
that an incorrect $\hat{X}_i$ (one with $\hat{X}_i \neq X_i$) is
used to update the predictor context or repair-stream parsing for
subblock $i + 1$ is at most $\delta$.}

\emph{\textbf{Proof.}} The framing assumption guarantees that
subblock $i + 1$'s repair-stream bits are isolated by the header
offset and not by any function of $\hat{X}_i$. The context update
rule for the predictor is conditional on $H_v(\hat{X}_i) = v_i$. If
$\hat{X}_i \neq X_i$, then $H_v(\hat{X}_i) = v_i$ can occur only as
a hash collision with the unique stored $v_i$. Under the
cryptographic hash assumption, this happens with probability at
most $\delta$ per candidate; the bound holds by union over the
considered candidates if $\delta$ is interpreted accordingly.

\emph{Caveat.} Without the explicit-framing assumption, a deletion
or length-mismatch in subblock $i$'s repair-stream encoding (e.g.,
caused by a non-self-delimiting per-mode encoding without a header
offset) shifts the byte boundaries of every subsequent subblock and
can make all later subblocks unreadable; the per-subblock hash $v_i$
detects that a failure occurred but does not localize or recover
later subblocks. The framing assumption (explicit per-subblock
offset table in the archive header, enforced by the engineering
specification) eliminates this failure mode.
\hfill$\blacksquare$

The theorem is stated for cryptographic $H_v$ because
non-cryptographic hashes (CRCs, simple checksums) admit structured
collisions that an adversarial or simply unlucky predictor may
produce. In practice we recommend a truncated SHA-256 per verified
subblock. A 64-bit truncation gives $\delta \approx 2^{-64}$, which
is operationally negligible and costs 8 bytes per subblock; a full
SHA-256 costs 32 bytes.

\textbf{10. Practical considerations}

\textbf{10.1 Bit-exact decoder reproduction of $M$}

All RNR types require that the decoder reproduce the model output
$Y = M(C), Q,$ or $U$ bit-for-bit identically to the encoder. IEEE-754
floating-point arithmetic is non-associative, GPU reduction kernels
are non-deterministic by default, and cuDNN/cuBLAS algorithm selection
varies by hardware and library version. A model executed naively on a
heterogeneous fleet of decoders will not produce identical output, and
a single bit of divergence breaks lossless decoding.

Three remedies are available.

\begin{itemize}
\item
  Integer arithmetic. Quantize the model to 8‑bit or 16‑bit integers and
  execute inference in pure integer kernels. This is what NNCP uses in
  its production CPU pipeline. Integer arithmetic is associative and
  hardware‑independent up to the well‑defined behavior of the integer
  operations.
\item
  Fixed‑point arithmetic. Use a fixed‑point representation with explicit
  scale parameters and rounding modes. This is more flexible than pure
  integer arithmetic and is the standard approach in deterministic
  embedded inference.
\item
  Strict‑mode floating‑point. Use the strict IEEE‑754 mode of the
  inference runtime, disable fast‑math compiler optimizations, force a
  fixed reduction order, and pin the library version. This is fragile
  but sometimes acceptable for distribution alongside a specific runtime
  binary.
\end{itemize}

We recommend integer arithmetic with explicit per-layer scale and
zero-point parameters as the default for any RNR realization. The
model description $L(M)$ then includes both the parameter table and
the per-layer quantization scheme.

\textbf{10.2 Model amortization regimes}

$L(M)$ is paid once per archive. For a single-block archive, the model
cost dominates and RNR is competitive only if the model is small (a
few kilobytes), or if the model is pre-shared between encoder and
decoder via a public release. For a multi-block archive, the amortized
cost is $L(M)/m$ where $m$ is the number of blocks; for streaming
applications where $m$ is large, $L(M)$ effectively vanishes. The
three regimes have different optimal model sizes and warrant separate
evaluation.

\textbf{10.3 Verification overhead analysis}

A SHA-256 per 4 KiB subblock is $32/4096 = 0.78\%$ overhead. A
truncated 64-bit hash per 4 KiB subblock is $8/4096 = 0.20\%$. A
truncated 32-bit hash is $0.10\%$, but with a $2^{-32}$ collision
probability that may be acceptable in some applications. Reducing the
subblock size to 1 KiB increases overhead proportionally. The right
balance depends on the underlying error rate of the predictor; if
errors are rare, larger subblocks are acceptable.

\textbf{10.4 Fallback modes}

Every RNR coder must include a fallback mode that handles the case in
which the chosen primary mode fails to compress. Two fallbacks are
standard: (i)~raw storage of the subblock, costing $8N$ bits plus
headers; (ii)~baseline coder storage, applying a fixed external coder
(e.g., zstd) to the subblock. The choice between fallbacks is
per-subblock and recorded in the mode-flag stream of the Type-III-A
scheduler. The cost of the fallback flag must be included in
$L(\Theta)$.

\textbf{10.5 A sufficient condition for bit-exact reproduction}

Section 10.1 recommends integer arithmetic as the default determinism
mechanism but does not formalize when it suffices. The following theorem
gives a sufficient condition under which a quantized neural predictor is
bit-exact across any pair of conforming implementations.

\textbf{Theorem 10.1 (Bit-exact reproduction under deterministic
integer inference).} \emph{Let $M$ be an inference graph with $L$
layers, where layer $\ell$ is implemented as an integer matrix-vector
product of fan-in $d_\ell$ followed by a non-linearity. Suppose:
(i)~all weights and inputs are integers; (ii)~every accumulator in
layer $\ell$ uses signed bit-width $b_\ell$ satisfying
\[
b_\ell \geq \lceil \log_2(d_\ell \cdot W_\ell \cdot I_{\ell-1} + 1) \rceil + 1,
\]
where $W_\ell$ is an upper bound on the magnitude of any weight entry
at layer $\ell$ and $I_{\ell-1}$ is an upper bound on the magnitude of
any input to layer $\ell$ (the $+1$ inside the $\log$ accounts for
the asymmetric signed-two's-complement range $[-2^{b-1}, 2^{b-1}-1]$,
so the positive extremum $+d_\ell W_\ell I_{\ell-1}$ is representable); (iii)~all non-linearities are deterministic
integer-to-integer mappings (e.g., piecewise-constant lookup tables);
(iv)~all reduction operations (sums in matrix-vector products) follow
a fixed, deterministic order specified by the model description. Then
no accumulator overflows and $M(C)$ produces bit-identical output for
any input $C$ on any two implementations conforming to (i) through
(iv).}

\emph{\textbf{Proof.}} Under (i) and (ii), each integer multiplication
and addition produces a result within the representable range; no
overflow or saturation occurs. Integer arithmetic is associative and
commutative, so any ordering of additions produces the same exact sum.
Under (iv), the ordering is pinned by the model description, so two
implementations choose the same ordering and produce the same sum.
Under (iii), the non-linearity is a deterministic function from
integer inputs to integer outputs that is fully specified by the model
description (typically as a lookup table); two implementations
applying the same table to the same input produce the same output. By
induction on layer depth, the output of layer $l$ is identical across
implementations, hence the final $M(C)$ is identical.
\hfill$\blacksquare$

The bit-width condition is per layer. If each layer renormalizes or
clips outputs so that $|x| \leq I_{\max}$ holds uniformly, then a
uniform sufficient condition is
$b \geq \lceil \log_2(d_{\max} \cdot W_{\max} \cdot I_{\max} + 1) \rceil + 1$,
where $d_{\max} = \max_\ell d_\ell$ is the maximum fan-in (the $+1$
inside the $\log$ matches the asymmetric signed-integer range, as in
the per-layer bound). Without inter-layer renormalization,
intermediate magnitudes can grow recursively:
$I_\ell \leq d_\ell \cdot W_\ell \cdot I_{\ell-1}$ in the worst
case, so $I_\ell \leq I_0 \prod_{r \leq \ell} d_r W_r$ and depth
enters multiplicatively, not linearly. Numerical sanity check: for
int8 quantization with $I_{\max} = W_{\max} = 2^7$ and fan-in
$d_{\max} = 2^{12}$ (4096-channel matrix-vector product), the
worst-case sum is bounded by $2^{12} \cdot 2^7 \cdot 2^7 = 2^{26}$,
so $b = 32$ suffices with comfortable headroom. For int16
quantization with $I_{\max} = W_{\max} = 2^{15}$ and the same
fan-in, the worst-case sum is $2^{12} \cdot 2^{15} \cdot 2^{15} =
2^{42}$, so $b \geq 44$ is needed (in practice $b = 64$ is the next
natural width; $b = 32$ overflows by 10+ bits at this fan-in;
only the degenerate scalar case $d_{\max} = 1$ gives a worst-case
product of at most $2^{30}$ and fits $b = 32$; any matrix-vector
layer with fan-in $d_{\max} \geq 2$ under int16 weight and
activation bounds exceeds $2^{31}$ in the worst case and requires
$b \geq 33$). int16-only deployments must
either use 64-bit accumulators or insert per-layer renormalization
to bound $I_\ell$ at runtime. The non-linearity table requirement
excludes any function involving real-valued division, square root, or
transcendentals; sigmoid, tanh, GELU, and softmax must be replaced by
piecewise-linear or integer-quantized approximations, which is
standard practice in quantized inference. The reduction ordering
requirement excludes any platform-specific parallelization (e.g., GPU
atomics) that does not commit to a deterministic order; this is
satisfied by sequential CPU implementations and by GPU implementations
that use tree-reduction with a fixed schedule.

Theorem 10.1 reduces the reproducibility question from a software
engineering problem to a precision-and-determinism specification. A
reference implementation that meets the four conditions is bit-exact
reproducible across any platform that conforms; cross-platform
conformance becomes a testable property rather than an empirical hope.
Constructing such an implementation and certifying conformance is
engineering work and is deferred to a companion artifact (§13.3).

\textbf{10.6 Bit-exact adaptive predictors under integer-arithmetic
optimization}

Theorem 10.1 covers the forward pass of a fixed predictor $M$. For RNR
coders that adapt $M$ online from previously verified subblocks during
encoding --- a useful capability for streams whose statistics drift
over a long archive --- the adaptation must itself be deterministic
and decoder-reproducible. Otherwise the encoder and decoder would
diverge after the first adaptation step. The relevant discipline is
integer-arithmetic optimization, developed in the recent
low-precision-training literature.

\textbf{Definition 10.1 (Compatible online optimizer).} An online
optimizer $\mathrm{Opt}$ is compatible with bit-exact reproduction if:
(a)~its state at step $t$, denoted $S_t$ (parameters, momentum,
accumulators, and any auxiliary buffers), is a tuple of
fixed-bit-width integers; (b)~its update function
$S_{t+1} = \mathrm{Opt}(S_t, g_t)$ takes the current state and a
fixed-bit-width integer gradient $g_t$ and produces the new state via
a deterministic computation using only integer arithmetic, lookup
tables, and seeded pseudorandom number generation; (c)~any operation
that would otherwise require real-valued computation (division, square
root, exponential, transcendentals) is replaced with a specified
deterministic integer approximation (integer division with explicit
rounding mode, integer square root, fixed-point or lookup-table-based
approximation of transcendentals).

Several practical optimizers fit Definition 10.1. Integer-quantized
SGD with quantized gradients (Wu, Yan, Wang, Maxson, Xu, He 2018) is
the canonical example: parameters, gradients, and
learning-rate-scaled updates are all integers, and the update is
$\text{parameter} \mathrel{-}= \mathrm{shift\_right}(\text{gradient} \cdot \text{learning\_rate}, \text{scale})$.
SignSGD (Bernstein, Wang, Azizzadenesheli, Anandkumar 2018) is
compatible by construction since the only non-trivial operation is
sign extraction. Integer-arithmetic Adam (Banner, Hubara, Hoffer,
Soudry 2018) is compatible when the reciprocal-square-root operation
is replaced with a deterministic integer approximation (typically a
piecewise-linear lookup table). Non-examples include any
floating-point optimizer without quantization wrappers, Adam with
real-valued square roots, and stochastic optimizers with unseeded
random masking.

\textbf{Theorem 10.2 (Bit-exact adaptive predictors, causal updates).}
\emph{Suppose the forward pass of $M$ satisfies the conditions of
Theorem 10.1, and the \emph{backward pass} (integer backpropagation)
satisfies the same four conditions independently: (i)~all backward
operands and gradient buffers are integers, (ii)~every accumulator
in any backward matrix-vector or reduction has signed bit-width
$b^{\mathrm{bwd}}_\ell$ satisfying the per-layer bound of
Theorem 10.1 applied to the backward operation's fan-in, weight
magnitude, and incoming-gradient magnitude bounds; (iii)~every
derivative-of-non-linearity is a deterministic integer-to-integer
lookup table; (iv)~every backward reduction follows a fixed,
deterministic order specified by the model description. Suppose the
encoder updates $M$ online via a compatible online optimizer
(Definition 10.1) using the gradient computed by this backward
pass. Suppose the order in which subblocks are used for adaptation is
pinned by the model description and identical on encoder and
decoder. Suppose further the \emph{causal-update condition}: the
encoder's adaptation step at subblock boundary $t$ uses only data
already available to the decoder at that boundary (i.e., bytes from
subblocks $0, \ldots, t$ all of which the decoder has already
verified per Theorem 9.1 before predicting subblock $t+1$). Then at
every subblock boundary, the encoder's and decoder's parameter
states are bit-identical, and the resulting predictor outputs are
bit-identical. \emph{Caveat:} without the causal-update condition
(e.g., if the encoder updates using the current byte before encoding
it), the decoder cannot replicate the same update sequence, and
parameter states diverge despite deterministic backprop and pinned
order; without the backward-pass conformance conditions, two
implementations sharing a compatible optimizer can still compute
different gradients (e.g., one wrapping, another saturating an
accumulator overflow) and so diverge.}

\emph{\textbf{Proof.}} Induction on the adaptation step $t$. Base
case ($t = 0$): the initial parameter state is part of $L(M)$ and is
bit-identical by construction. Induction step ($t \to t+1$): assume
$S_t$ is bit-identical between encoder and decoder. The subblock
$X_t$ used for adaptation is verified by Theorem 9.1 to be the same
on both sides. The gradient $g_t = \nabla \mathcal{L}(X_t; S_t)$ is
computed by integer backpropagation satisfying the four backward
conformance conditions (i)--(iv) above; by the same argument as in
Theorem 10.1 (no overflow, deterministic reduction order,
deterministic non-linearity-derivative tables, integer inputs),
each backward layer's output is bit-identical across implementations,
and by induction on backward depth so is the full gradient
$g_t$. The update $S_{t+1} = \mathrm{Opt}(S_t, g_t)$ applies
the deterministic integer-to-integer function from Definition 10.1
to bit-identical inputs, producing bit-identical outputs. By
induction, $S_t$ is bit-identical for all $t$. By Theorem 10.1
applied to each $S_t$, the predictor output is bit-identical at every
adaptation step. \hfill$\blacksquare$

Theorem 10.2 reduces the adaptive-predictor reproducibility question to
the existence of a compatible online optimizer with adequate training
performance. The low-precision-training literature provides several such
optimizers with published convergence and accuracy results approaching
their floating-point counterparts. The remaining engineering work is to
wire one of them into a concrete RNR codec; this is deferred to a
companion artifact (§13.3) on the same grounds as Theorem
10.1\textquotesingle s reference implementation.

\textbf{10.7 Random access}

Theorems 10.1 and 10.2 concern bit-exact decoder reproduction. A
separate, operationally important question is whether a decoder can
extract a small portion of $X$ without decoding all of it --- the
seek-and-read functionality of bgzip and zstd seekable mode.
Existing neural compressors lack this capability for two distinct
reasons: autoregressive byte/pixel models (NNCP, PixelCNN++,
PixelSNAIL) require sequential evaluation of all $p$ previous bytes
before byte $p$ is available (cost $\Theta(p)$ per random read);
invertible-flow and bits-back block codecs (IDF++, Bit-Swap,
HiLLoC) operate on whole blocks of statically determined size and
state, so even a one-byte read forces decoding of the enclosing
block and any state on which it depends. In either case the
effective minimum random-read granularity is large and growing with
archive size (the first case linearly in $p$; the second
block-by-block with no sub-block seek). This is a deployment
blocker for archive-style use cases where partial reads are common.
The RNR framework supports random access in time independent of
$|X|$, provided the predictor has bounded context. We make this
precise.

\textbf{Definition 10.2 ($W$-bounded predictor).} A predictor $M$ is
$W$-bounded if its output $Y_i = M(C_i)$ at position $i$ depends only
on the previous $W$ bytes of public context:
$Y_i = M(X_{i-W}, \ldots, X_{i-1})$ when $C$ is built incrementally
from decoded subblocks. The bound $W$ is a structural property of
$M$'s architecture, not a tunable parameter at inference time.

Transformers with fixed context window $W$ are $W$-bounded, as are
$n$-gram and finite-state predictors of order $W$ and any feedforward
architecture with bounded receptive field. RNNs, LSTMs, GRUs, and
state-space models (Mamba, S4) are not $W$-bounded in general; their
state can depend on arbitrarily long history. Theorem 10.4 below
covers the non-$W$-bounded case via state snapshots.

\textbf{Definition 10.3 (Sync point).} A sync point at position $P$ in
the source is a position where (a)~the predictor's context state is
reset to a known initial value $Q_{\text{init}}$ independent of bytes
$[0, P)$, and (b)~the entropy coder's internal state (arithmetic-coder
interval, ANS register, or analogous) is flushed at the encoder and
reset to a known initial value at the decoder, so that the next bit
in the repair stream begins a fresh entropy-coder block. The encoder
inserts sync points at fixed intervals $K$ so that the byte offset
in $X$ of the $j$-th sync point is the deterministic value $j K$
(no byte-offset storage required), and includes a seek index in the
archive mapping each sync point identifier to its bit offset in $R$
only. The seek index has size at most
$\lceil |X|/K \rceil \cdot \lceil \log_2 L(R) \rceil$ bits.
(Non-uniform sync-point spacing would require storing $\log_2 |X|$
additional bits per entry, raising the seek index size to
$\lceil |X|/K \rceil \cdot (\lceil \log_2 L(R) \rceil + \lceil
\log_2 |X| \rceil)$ bits; we assume uniform spacing throughout.)
The entropy-coder flush adds $O(1)$ bits of termination per sync
point and contributes $O(|X|/K)$ bits of additional repair-stream
overhead in total; the predictor and entropy-coder reset together
make each sync-point bit offset a self-contained restart point for
decoding.

Each sync point forfeits some compression efficiency because the
predictor restarts cold and cannot exploit the prior $W$ bytes of
context. \emph{Bounded-cost assumption.} We assume the predictor's
output distribution satisfies a probability floor $\eta > 0$: for
every byte $b \in \{0, \dots, 255\}$ and every context $c$,
$\Pr_M(b \mid c) \geq \eta$ (in practice enforced by mixing with a
small uniform escape, $\eta \geq 2^{-12}$ being standard). Under
this floor, the per-byte cold-context arithmetic-code length is
bounded above by $\log_2(1/\eta)$. The per-sync-point excess cost
over warm-context encoding of the same $W$ bytes is therefore
bounded by
\[
\Delta_{\mathrm{sync}} \leq \sum_{i=P}^{P+W-1}
\bigl[\log_2(1/\eta) - (-\log_2 \Pr_M(X_i \mid C_i^{\mathrm{warm}}))\bigr]
\leq W \cdot \log_2(1/\eta),
\]
where $C_i^{\mathrm{warm}}$ denotes the warm $W$-byte context the
encoder would have used in a sync-point-free encoding. Per-byte
cold cost is at most $\log_2(1/\eta)$ and warm cost is at least
$0$, so each summand is at most $\log_2(1/\eta)$ and at least
$-\log_2(1/\eta)$ (the latter when warm cost exceeds the cold-cost
upper bound, which can happen on very unpredictable data); the
trivial upper bound $W \cdot \log_2(1/\eta)$ is loose and does not
depend on the warm rate. Sharper upper bounds require additional
assumptions: e.g., a uniform per-byte warm-context \emph{lower}
bound $r^{\mathrm{warm}}_{\min}$ on $-\log_2 \Pr_M(X_i \mid
C_i^{\mathrm{warm}})$ valid at all positions, yielding
$\Delta_{\mathrm{sync}} \leq W \cdot (\log_2(1/\eta) - r^{\mathrm{warm}}_{\min})$.
(An upper bound on warm cost would give a lower, not an upper,
bound on $\Delta_{\mathrm{sync}}$, so the warm-rate assumption must
be a lower bound to yield a useful overhead ceiling.) Without the
probability-floor assumption (e.g., a predictor that can place
arbitrarily small mass on the true byte after reset), the
per-sync-point cost is unbounded. The total compression cost of
sync points under the floor assumption is at most
$(|X|/K) \cdot W \cdot \log_2(1/\eta)$ bits.
Concrete-number estimates in the abstract and §11 (e.g., the
``approximately one percent'' figure) are heuristic conjectures
based on empirical observations of cold-vs-warm rates for small
distilled transformers; they are not theorems and depend on the
sharper assumption above.

\textbf{Theorem 10.3 (Random access, output-sensitive).} \emph{Let
$X$ be compressed under RNR with a $W$-bounded predictor $M$
satisfying Theorem 10.1 and with sync points placed at uniformly
spaced source positions $j K$ for $j = 0, 1, \dots$ (Definition
10.3). Let $|X|$ denote the archive's uncompressed source length.
Random access to $k$ consecutive bytes starting at any position
$p \in [0, |X|)$ requires: (i)~one direct seek-index lookup at
entry $\lfloor p/K \rfloor$, reading $O(\log L(R))$ bits to obtain
the repair-stream bit offset of the sync point at byte position
$P = K \lfloor p/K \rfloor$; (ii)~decoding of at most $K + k$
positions of $R$, of which at most $K$ positions are warmup
(discarded output); (iii)~$K + k$ predictor evaluations and
entropy-coder updates. Total work is
$O(\log L(R) + (K + k) \cdot \mathrm{cost}(M))$. (For non-uniform
sync-point placement, the lookup becomes a binary search costing
$O(\log(|X|/K) \cdot \log L(R))$ bits; the rest of the analysis is
unchanged.) \emph{This cost is polylogarithmic in archive size
$|X|$ only when $K$ and $k$ are held fixed or grow at most
polylogarithmically in $|X|$ and $L(R) = \mathrm{poly}(|X|)$; for
$K = \Theta(|X|)$ or $k = \Theta(|X|)$ the cost is
$\Theta(|X| \cdot \mathrm{cost}(M))$, no better than full sequential
decode}. When $(K+k) \cdot
\mathrm{cost}(M) \gg \log L(R)$ the second term dominates and the
cost is $O((K+k) \cdot \mathrm{cost}(M))$, independent of $|X|$ in
that parameter regime.}

\emph{\textbf{Proof.}} Under the uniform-spacing convention of
Definition 10.3, the byte position of the latest sync point
$P \leq p$ is the deterministic value $P = K \cdot \lfloor p/K \rfloor$
(a constant-time computation, no search required). Look up the
corresponding bit offset in $R$ from the seek index: the index has
exactly $\lceil |X|/K \rceil$ entries, each of width
$\lceil \log_2 L(R)\rceil + O(1)$ bits, indexed by $\lfloor p/K \rfloor$
in $O(\log L(R))$ bits per direct probe. (For non-uniform spacing,
binary search over the index would cost
$O(\log(|X|/K) \cdot \log L(R))$ bits in total; the uniform case
avoids the $\log(|X|/K)$ factor at the cost of one extra
constant-time arithmetic operation.) By
Definition 10.3 (a, b), decoding from $R$'s bit offset at $P$
proceeds with both the predictor state initialized to
$Q_{\text{init}}$ and the entropy-coder state initialized to its
post-flush starting value; no prior bits of $R$ are required to
parse the bits following the sync-point offset. Decode forward one
position at a time: at each position $i \in \{P, P+1, \ldots, p+k-1\}$,
compute $Y_i = M(C_i)$ where $C_i$ is the context built from
previously decoded positions $[P, i)$. By the $W$-bounded property
(Definition 10.2), $Y_i$ depends only on bytes $[i-W, i)$, which are
all available within the decoded prefix once $i$ exceeds $P + W$. The
first $W$ positions after $P$ may have shortened context (positions
$[P, P+W)$ see fewer than $W$ bytes), which the encoder accommodates
by encoding those positions with reduced-context predictions; this is
exactly the cold-context cost analyzed before. Continue decoding
until position $p + k$; output bytes $[p, p+k)$. By Theorem 10.1, all
predictor evaluations are deterministic and bit-identical to a full
decoding from position 0, so the output matches the original $X$ on
$[p, p+k)$. Total decoded positions: $p + k - P \leq K + k$ since
$P > p - K$. \hfill$\blacksquare$

Theorem 10.3 establishes that random access in RNR archives runs in
time with only polylogarithmic dependence on archive size (a
$\log(|X|/K) \cdot \log L(R)$ seek-index cost plus a $K+k$ decode
window). For $K = 64$ KiB and $k = 1$ KiB, the cost is roughly
$K + k = 65$ KiB worth of predictor evaluations --- typically on
the order of $650\,\text{ms}$ on a single CPU core for a small
distilled transformer (one forward pass at $\sim 10\,\mu\text{s}$
per byte; faster on GPU/NPU). This is asymptotically a
factor of $|X|/(K + k)$ faster than full sequential decoding, which for
archive sizes of 1 GiB and $K = 64$ KiB is a roughly 16{,}000-fold
speedup.

\textbf{Theorem 10.4 (Random access for non-$W$-bounded predictors).}
\emph{For a non-$W$-bounded predictor (RNN, LSTM, state-space model),
the predictor-reset condition (a) of Definition 10.3 cannot be
satisfied by cold-resetting context alone: the predictor's
recurrent hidden state depends on arbitrarily long history. Replace
condition (a) by an \emph{augmented sync point}: instead of
resetting the predictor's hidden state, store its full state at
position $P$ in the seek-index entry (so the decoder, when seeking
to $P$, loads the stored state and proceeds with bit-exact
continuation). Condition (b) on entropy-coder flush/reset remains
unchanged. The augmented snapshot stores the predictor's recurrent
hidden state, any adaptive parameter snapshot (if Theorem 10.2 is
in effect), and the entropy-coder boundary state (range/state for
arithmetic coding or ANS state, if a different flush convention is
used). Let $L(\text{state})$ be the bit length of this combined
boundary snapshot. Under the uniform-spacing convention of
Definition 10.3, the augmented seek index has exactly
$\lceil |X|/K \rceil$ entries (one per uniform sync point), with
total size $O((|X|/K) \cdot (\log L(R) + L(\text{state})))$ bits.
The random-access cost per query (to $k$ bytes at position $p$) is
$O\bigl(\log L(R) + L(\text{state}) + (K + k) \cdot
\mathrm{cost}(M)\bigr)$ under uniform spacing (direct index lookup),
or $O\bigl(\log(|X|/K) \cdot (\log L(R) + L(\text{state})) + (K + k)
\cdot \mathrm{cost}(M)\bigr)$ under non-uniform spacing (binary
search), where the index-related terms cover the entry lookup plus
loading and deserializing the snapshot before decoding can
resume.}

\emph{\textbf{Proof.}} Augment each sync point's index entry with
all boundary state at position $P$, serialized as a fixed-bit-width
integer tuple as required by Theorem 10.1. The combined snapshot adds
$L(\text{state})$ bits per sync point, multiplying the index size by
$(1 + L(\text{state}) / \log L(R))$. Under uniform sync-point
spacing (Definition 10.3), to random-access bytes $[p, p+k)$, the
latest sync point $P \leq p$ is the deterministic value
$P = K \lfloor p/K \rfloor$ (constant-time computation); direct
lookup at entry $\lfloor p/K \rfloor$ reads
$O(\log L(R) + L(\text{state}))$ bits to obtain the bit offset and
snapshot. Under non-uniform spacing, the lookup is a binary search
over the augmented index, costing
$O(\log(|X|/K) \cdot (\log L(R) + L(\text{state})))$ bits in total.
In either case, deserialize the boundary snapshot, restore the
decoder's predictor, adaptive-parameter, and entropy-coder boundary
states, and decode forward as in Theorem 10.3 (cost
$O((K+k) \cdot \mathrm{cost}(M))$).
The decoded output is bit-identical by Theorem 10.1 applied to the
restored state. \hfill$\blacksquare$

For LSTM predictors with $d$-dimensional hidden state and
$d$-dimensional cell state (the standard LSTM has both), the
snapshot is $L(\text{state}) = 2 d \cdot b$ bits where $b$ is the
per-component bit-width. For modern state-space models with
structured state-transition matrices, $L(\text{state})$ can be
$O(d^2)$. With $d = 1024$ and $b = 16$, an LSTM state snapshot is
$2 \cdot 1024 \cdot 16 = 32{,}768$ bits $= 4$ KiB; placing one
snapshot per $K = 64$ KiB sync point adds approximately $6.25\%$
index overhead. State-space models
with $d = 256$ and structured-state representation can require 64 KiB
snapshots, doubling the seek index. This motivates the $W$-bounded
design choice for deployable RNR codecs, particularly for archive
formats where the seek index is co-located with the data and
contributes to its size.

\emph{\textbf{Trade-offs}}

Three trade-offs govern random-access design in RNR archives.

\begin{itemize}
\item
  Sync-point spacing $K$. Smaller $K$ means faster seek and smaller
  warmup cost, but more sync points and therefore more cold-context
  compression overhead. Typical choices: $K = 4$ KiB (aggressive
  seek), $K = 64$ KiB (balanced), $K = 1$ MiB (compression-favored).
\item
  Type-II semi-non-factorized coding. Theorem 5.4's $\Theta(N)$
  separation requires joint structure across all $N$ positions of a
  block. Inserting sync points every $K$ positions partitions the
  block into independent $K$-length subblocks, recovering only
  $\Theta(K)$ separation per subblock. For $K = 64$ KiB and $N = 1$
  GiB, the per-subblock advantage is preserved but the cross-subblock
  joint structure is forfeited. Type-I and Type-III-B are not
  affected by this trade-off because their constructions are local
  (per-position or per-tile within sync-bounded windows). Type-III-C
  is affected when dictionary tokens span longer fragments or
  grammar expansions. A token that would cross a uniform Def-10.3
  sync-point boundary $jK$ must be split, with the affected
  occurrence forfeiting its dictionary saving. Over $|X|/K$ sync
  points the worst-case loss is bounded by
  $|X| \cdot \max_\tau |\tau| / K$ bits (at most one token crosses
  any given boundary, and the lost saving per crossed token is at
  most $\max_\tau |\tau|$ bits), negligible for $K \gg \max_\tau
  |\tau|$ but a real cost when $K$ is small (e.g., $K = 4$ KiB and
  grammar tokens of $\sim 1$ KiB give a few per-cent overhead). A
  non-uniform sync-point variant that aligns boundaries to token
  edges instead is also available, at the cost of storing byte
  offsets in the seek index (Definition 10.3's non-uniform spacing
  case, adding $\lceil \log_2 |X|\rceil$ bits per entry); the
  random-access cost then follows the non-uniform branch of
  Theorem 10.3 ($O(\log(|X|/K) \cdot \log L(R))$ binary-search
  bits).
\item
  Verification under random access. Theorem 9.1's chained-hash
  verification covers sequential decoding; a random-access reader
  cannot verify bytes $[0, p)$ it never decoded. Three solutions
  exist: (a)~per-sync-point hashes that allow local verification,
  costing $O((|X|/K) \cdot |\text{hash}|)$ bits; (b)~a Merkle tree
  over the repair stream allowing verified random access with
  $O(\log |X|)$ hash overhead per read; (c)~explicitly unverified random access for
  trusted-source archives. The Merkle-tree option is standard in
  content-addressable storage and is the recommended default for
  archive deployments.
\end{itemize}

\emph{\textbf{Comparison with existing systems}}

Standard random-access lossless compressors:

\begin{itemize}
\item
  bgzip (block-gzip): independent gzip blocks, ratio cost
  approximately 3 to 7 percent versus monolithic gzip, random access
  in $O(K + k)$ time. Widely used in bioinformatics for indexed FASTA
  and BAM files.
\item
  zstd seekable mode: independent zstd frames with a footer index,
  ratio cost approximately 2 to 10 percent versus monolithic zstd,
  random access in $O(K + k)$ time.
\item
  LZMA seekable wrappers (ad hoc xz multi-stream archives): ratio
  cost approximately 5 to 15 percent, random access in $O(K + k)$
  time.
\end{itemize}

Existing neural lossless compressors lack sub-archive random access,
with two distinct underlying obstructions. Autoregressive byte and
pixel models (NNCP, PixelCNN++, PixelSNAIL) require sequential
evaluation of bytes $[0, p)$ before byte $p$ can be decoded; cost is
$\Theta(p)$ predictor evaluations. Block-codec bits-back and
invertible-flow models (Bit-Swap, IDF++, HiLLoC) decode in fixed
block/patch units, so even a one-byte read forces decoding of the
enclosing block; effective seek granularity is the block size, not
the byte (typical block sizes are kilobytes to megabytes). Neither
family supports byte-resolution random access as currently designed.
For a 1 GiB archive the autoregressive case is hours of wall-clock
time per access; the block-codec case decodes the enclosing block
on every read.

RNR with a $W$-bounded transformer predictor and sync points every
$K$ positions: random access in $O(K + k)$ predictor evaluations
independent of $|X|$, at a per-sync-point worst-case excess of
$W \cdot \log_2(1/\eta)$ bits (Section 10.7), giving a worst-case
total of $(|X|/K) \cdot W \cdot \log_2(1/\eta)$ bits over the
archive. For $K = 64$ KiB, $W = 4$ KiB, $\eta = 2^{-8}$, and a 1
GiB source, the worst-case bound evaluates to
$16384 \cdot 4096 \cdot 8$ bits $= 64$ MiB, or roughly 6.25\% of
the source size. The realistic overhead is smaller because the
per-byte excess in practice is the cold-vs-warm \emph{rate gap}
$\delta$ (typically a few tenths of a bit per byte for natural
data) rather than the full $\log_2(1/\eta) = 8$ bits. With
$\delta \approx 0.5$ bit/byte, $(|X|/K) \cdot W \cdot \delta =
16384 \cdot 4096 \cdot 0.5$ bits $\approx 32$ Mibits $\approx 4$
MiB, or roughly 0.4\% of the source size; this estimate is
heuristic and is not a theorem. The ratio comparison
against bgzip and zstd seekable (typically 2--7 percent overhead)
is empirical and is the subject of §11.

This is the framework\textquotesingle s deployment-grade distinguishing
feature relative to other neural compressors. The bounded-context
architecture, the deterministic reproducibility guarantee (Theorem
10.1), and the sync-point construction (Definition 10.3) jointly
enable seek-and-read access at a small constant fractional ratio
overhead --- a capability that LSTM-based or state-space-based
neural compressors can only provide with the expensive state
snapshots of Theorem 10.4 (4 KiB or more per sync point); that
NNCP-style autoregressive transformer compressors lack only because
their encoders do not insert sync points or maintain a seek index;
and that bits-back / invertible-flow block codecs (Bit-Swap, IDF++,
HiLLoC) lack because their block/ANS/latent decoding requires
recovering the enclosing block (or vectorized batch of patches) as
a unit, so the effective seek granularity is the block, not the
byte; sub-block-granularity seek is not provided by the current
designs without modifying their block structure.

\textbf{10.8 Block-device and filesystem deployment}

A natural application of the random-access result is to expose an
RNR-compressed archive as a read-only block device or filesystem through
an operating-system driver, allowing files or sectors to be accessed on
demand without ever decompressing the whole archive. This is the
deployment pattern of squashfs, EROFS, btrfs with zstd compression, and
Apple Disk Images: read-only compressed system images mounted by the
kernel and served to user-space applications transparently. Existing
neural compressors cannot serve this use case because their
$O(|X|)$ sequential decoding requirement makes any access to byte
$p$ cost the same as decompressing all preceding bytes; for archive
sizes typical of disk images (gigabytes to terabytes), this is
prohibitive. RNR removes this obstruction. The two theorems below
formalize the application.

\textbf{Theorem 10.5 (Block-device interface, parameter-conditional).}
\emph{Let $X$ be compressed under RNR with a $W$-bounded predictor
$M$ satisfying Theorem 10.1 and sync points placed every $K$
positions; let $|X|$ denote the uncompressed source length. A
block-device driver that maps logical-block-address-indexed read
requests (LBA $\mathcal{L}$, length $k$ bytes) to RNR random-access
calls can serve any read request in time
$O(\log(|X|/K) \cdot \log L(R) + (K + k) \cdot \mathrm{cost}(M))$ and
working memory $O((K + k) \cdot 8 + L(M) + L(\mathrm{seek\_index}))$
\emph{bits}, equivalently $O(K + k + L(M)/8 + L(\mathrm{seek\_index})/8)$
\emph{bytes}, where $L(M)$ is the predictor description size in bits
(kept resident in the driver) and $L(\mathrm{seek\_index})$ is the
seek-index size in bits (typically $O((|X|/K) \log L(R))$). The cost is polylogarithmic in
archive size $|X|$ only in the parameter regime where $K$, $k$, and
$\log L(R)$ are held fixed or grow at most polylogarithmically in
$|X|$; outside that regime the output-sensitive cost of Theorem 10.3
applies.}

\emph{\textbf{Proof.}} Each LBA-indexed read request $(\mathcal{L},
k)$ corresponds to a byte-range request on the underlying source $X$
starting at byte position
$p = \mathcal{L} \cdot \text{sector\_size}$ and continuing for $k$
bytes. By Theorem 10.3, RNR random access to $k$ bytes at position
$p$ runs in $O(\log(|X|/K) \cdot \log L(R) + (K + k) \cdot \mathrm{cost}(M))$ time and
requires $O(K + k)$ bytes (equivalently $O((K+k) \cdot 8)$ bits) of
decode-window working memory (warmup buffer plus output buffer).
The driver additionally keeps the predictor weights ($L(M)$ bits)
and the seek index ($L(\mathrm{seek\_index})$ bits) resident for
constant-time per-query lookup; these are fixed at archive-mount
time, not per-query. The translation from $(\mathcal{L}, k)$ to
$(p, k)$ is a constant-time arithmetic operation. The total
per-request cost is therefore as stated, with the parameter regime
caveat inherited from Theorem 10.3. \hfill$\blacksquare$

The $\mathrm{cost}(M)$ factor governs whether the driver is fast
enough for interactive use. For a small distilled byte-level
transformer (10 to 50 million parameters, int8 weights), one
single-byte forward pass costs approximately $10\,\mu\text{s}$ on a
single CPU core, $1\,\mu\text{s}$ on a server GPU, and
$0.1\,\mu\text{s}$ on a high-end NPU (engineering estimates; actual
numbers vary by quantization scheme and hardware generation).
Treating one forward pass as one decoded byte, per-byte inference
throughput is thus approximately $100$ KB/s on CPU, $1$ MB/s on
GPU, and $10$ MB/s on NPU. For $K + k = 64$ KiB, the random-access
latency for cold reads is therefore approximately
$64 \text{ KiB} \cdot 10\,\mu\text{s/byte} \approx 650\,\text{ms}$
(CPU), $\approx 65\,\text{ms}$ (GPU), and
$\approx 6.5\,\text{ms}$ (NPU). The NPU end of this range is in the
same order of magnitude as a slow random-read of a multi-MB region
from local NVMe (single-block NVMe random-read latency is
$\sim 100\,\mu\text{s}$, so RNR-NPU is about $65 \times$ slower per
request, comparable to a server-side decompressor); the CPU end
matches spinning-disk seek latency. Hardware-batching multiple
forward passes per kernel invocation can improve effective
throughput by a factor of 10--100 on GPU/NPU when access patterns
are sequential, narrowing the cold-read latency proportionally.
Production deployments will target NPU or GPU paths when low latency
is required and fall back to CPU only for archive inspection or
compatibility.

\textbf{Theorem 10.6 (Hardware-target portability, explicit
conformance).} \emph{Let $M$ satisfy the integer-arithmetic,
lookup-table, and fixed-reduction-order conditions of Theorem 10.1.
A hardware target $T$ is said to \emph{conform} for $M$ if it pins,
in a target-specific test vector matched to the engineering spec,
all of the following: (a)~signed integer representation (two's
complement) at the bit-widths declared by the model description;
(b)~the result of every left/right shift, division, and modulo at
those bit-widths (including round-toward-zero vs floor and signed
vs unsigned semantics); (c)~the saturation-vs-wrap behaviour at each
arithmetic operation; (d)~the byte-level serialization of every
non-linearity lookup table; (e)~the reduction-tree shape and order;
(f)~the per-target serialization of all approximation functions
(integer Newton-Raphson, table-lookup interpolation, etc.) used in
the forward pass; (g)~the byte-level encoding of all model
parameters and intermediate tensors, including endianness, signed
vs unsigned interpretation of each declared bit-width, the packing
order of sub-byte quantizations (e.g., int4-pair-per-byte: low
nibble first vs high nibble first), zero-point and scale storage
layout, and tensor index-to-memory order for any layout-sensitive
tensor (e.g., NHWC vs NCHW for convolution weights; row-major vs
column-major for matrix multiplications). Under (a)-(g), the output
$M(C)$ computed on $T$ is bit-identical to that computed on any
other conforming target. The choice of target affects throughput
and latency but not the produced bytes; the engineering spec lists
the conformance tests that establish (a)-(g) for each named target
class.}

\emph{\textbf{Proof.}} Direct corollary of Theorem 10.1, with the
explicit observation that conformance is a property of the
target\textquotesingle s integer-arithmetic, lookup-table, and
data-layout semantics rather than of its parallelism strategy or
vendor identity. Two's-complement signed integer addition and
multiplication at a fixed bit-width are associative and commutative
\emph{within the representable range and under a pinned wrap-vs-saturate
convention}; condition (a) fixes the representation, (c) fixes the
wrap-vs-saturate behaviour for each operation, and the per-layer
bit-width bound of Theorem 10.1 keeps every accumulator within the
representable range so neither wrap nor saturation triggers in a
conforming forward pass --- under those conditions the produced
bits are identical across platforms. Conditions (b), (d), (e), (f),
and (g) pin the remaining sources of cross-platform divergence
(shift/division/modulo semantics, lookup-table layout, reduction
order, approximation functions, and weight/input byte
serialization); each is enforced by the engineering spec's
target-specific test vectors. Without (a)--(g), bit-identity is
not guaranteed: e.g., the same algorithm with different sub-byte
packing orders produces different weights. Fixed reduction order,
specified by the model description, is honored by any sequential
implementation and by any parallel implementation that uses a
tree-reduction with the specified schedule. Composing these properties
layer by layer through $M$ yields bit-exact output. \hfill$\blacksquare$

Conforming hardware targets include: scalar CPU implementations on any
architecture supporting the required integer bit-widths (x86-64, ARM64,
RISC-V, POWER); CPU SIMD implementations using integer SIMD instructions
(AVX2, AVX-512 with VNNI and IFMA extensions, ARM NEON, RISC-V vector
extension); GPU implementations using integer compute kernels and
integer Tensor Cores (NVIDIA Turing or later with int8 and int4 Tensor
Core paths, AMD CDNA with matrix accelerators); dedicated NPUs targeting
integer inference (Apple Neural Engine, Google Edge TPU, Qualcomm
Hexagon NPU, Huawei Ascend, AWS Inferentia, AWS Trainium in inference
mode). The same RNR archive can be encoded on a server GPU, served from
cloud storage, and decoded on a phone NPU, with byte-identity guaranteed
by Theorem 10.6 plus a conformance certification of each target.

\emph{\textbf{Cache amortization}}

A read-cache of size $C$ placed between the driver and the OS page
cache permits steady-state throughput exceeding the per-access bound
of Theorem 10.5. For sequential reads over any region of size up to
$C$, the amortized cost per byte is $O(\mathrm{cost}(M))$: one
predictor evaluation per byte, with no per-access warmup. For random
reads with spatial locality smaller than $C$, the cache hit rate
approaches one and the amortized cost per byte approaches
$O(\mathrm{memcpy})$ speed, dominated by main memory bandwidth. For
uniformly random reads at distance larger than $C$, the amortized
cost per byte equals the Theorem 10.5 bound divided by $k$, that is,
$O((K + k) \cdot \mathrm{cost}(M) / k)$, which approaches
$\mathrm{cost}(M)$ as $k$ grows. In typical deployments the
operating-system page cache provides an effective cache of several
gigabytes, absorbing read locality so that the worst-case Theorem
10.5 bound is the latency budget for cold reads only.

\emph{\textbf{Read-write semantics}}

RNR archives are read-only by construction: the predictor $M$ and the
sync-point structure assume a fixed source $X$. Three patterns
support read-write deployments.

\begin{itemize}
\item
  Copy-on-write overlay. The compressed archive is mounted as a
  read-only base; writes are directed to a separate writable layer
  (Linux overlayfs, unionfs, AUFS). Reads consult the overlay first
  and fall through to the base only on miss. This is the
  squashfs-plus-overlayfs pattern that container runtimes (Docker,
  Podman, containerd) and live Linux distributions use today.
\item
  Region rewrite. A modification within sync-point region
  $[P, P + K)$ re-encodes only that region: re-runs the predictor on
  the modified bytes starting from $P$, re-encodes the repair stream
  for the region, and updates the seek index entry for the affected
  sync point. Cost: $O(K \cdot \mathrm{cost}(M))$ per modification.
  The naive append-in-place form is well-defined only if the
  re-encoded region's bit length is bounded by a pre-reserved slot
  size; otherwise the rewritten region shifts every subsequent
  sync-point bit offset, and the archive needs either (i)~per-region
  fixed-size slots with reserved padding, (ii)~an indirection layer
  that decouples logical sync-point identifiers from physical bit
  offsets (e.g., a separately mutable offset table), or
  (iii)~rewriting all later seek-index entries on each write,
  raising the per-modification cost to
  $O(K \cdot \mathrm{cost}(M) + (N-P)/K)$ for the index update.
  Suitable for low-frequency writes on archives that benefit from
  in-place mutation, with the slot/indirection choice fixed at
  archive format design time.
\item
  Full re-encoding. Periodic offline compaction re-encodes the whole
  archive, merging accumulated overlay writes back into the base.
  Standard practice for read-mostly snapshot stores; the compaction
  cost is $O(N \cdot \mathrm{cost}(M))$ but is amortized over many
  writes.
\end{itemize}

For the disk-image use case named at the start of this section (compress
an ISO or block-device image, mount through a driver, use without full
decompression), the read-only base plus copy-on-write overlay pattern is
the natural deployment. It matches the architecture used today for
compressed container images and live system images, replacing the
dictionary-based compressed filesystem in the base layer with an RNR
archive of higher compression ratio at identical access semantics.

\emph{\textbf{Comparison with existing compressed filesystems}}

Compressed read-only filesystems are widely deployed:

\begin{itemize}
\item
  squashfs (Linux): gzip, xz, lzma, lz4, or zstd compression in 4 KiB to
  1 MiB blocks. Typical ratio 1.5 to 3 times. Random-access cost O(K +
  k). Used by live Linux distributions, snap packages, and container
  image layers.
\item
  EROFS (Linux): lz4, lzma, or deflate in 4 to 64 KiB blocks. Typical
  ratio 1.5 to 2.5 times. Random-access cost O(K + k). Used by Android
  system partitions and Huawei storage products.
\item
  btrfs with zstd (Linux): zstd compression in 4 to 128 KiB extents.
  Typical ratio 1.5 to 2 times. Random-access cost O(K + k).
\item
  APFS compressed files (macOS) and NTFS compressed files (Windows):
  LZ4-class compression at fixed sector boundaries. Typical ratio 1.2 to
  1.8 times. Random-access cost O(sector size).
\end{itemize}

RNR with a well-fitted transformer predictor is conjectured to
achieve compression ratios materially better than the dictionary-
based systems above on structured data classes (system images,
source archives, log files, structured binaries), at the same
random-access cost class and at the same conformance overhead.
Specific numerical ratios are not derived here and are not theorems;
order-of-magnitude predictions are stated as conjectures in §11
(``Type-I, Code12 on text and source code'', and the corresponding
predictions for other Types) and require empirical validation
against NNCP, PixelCNN++, LZMA, and zstd baselines. For
read-mostly use cases with rare writes --- container images, system
snapshots, archival storage, software distribution images --- the
qualitative deployment claim is that RNR matches the access pattern
of these systems while potentially offering a compression advantage
that empirical validation must establish. The cost is the predictor
inference compute, which is feasible at NPU or GPU speeds and acceptable
at CPU speeds for non-interactive workloads.

\emph{\textbf{Reference driver architecture}}

A reference RNR block-device driver consists of the following
components, mirroring the structure of existing compressed-filesystem
drivers:

\begin{itemize}
\item
  Header parser. Reads the archive header into memory: predictor
  weights, sync-point seek index, verification metadata, type
  metadata.
\item
  Predictor inference engine. Implements $M$ on one or more hardware
  paths (CPU scalar, CPU SIMD, GPU integer, NPU integer), each path
  conforming to Theorem 10.6. The path is selected at driver load
  time based on hardware availability.
\item
  Range-decode function. Given $(p, k)$, looks up the latest sync
  point $P \leq p$ in the seek index, decodes from $P$ to $p + k$
  using Theorem 10.3, returns bytes $[p, p + k)$.
\item
  Cache. LRU or ARC cache of recently decoded regions, with eviction
  policy tuned to expected access pattern. Sits between the
  range-decode function and the OS page cache.
\item
  OS interface. Translates kernel-level block-device or filesystem
  read requests into range-decode calls. Options include Linux FUSE
  for user-space deployment, Linux NBD or device-mapper for
  kernel-level deployment, the macOS FSKit interface, or a Windows
  file-system minifilter driver.
\end{itemize}

Every component except the predictor inference engine has direct
counterparts in existing compressed-filesystem drivers; the integration
work is to replace the dictionary-based decompressor with the neural
inference engine. Production implementations of small transformers on
each target hardware path already exist; combining them into a
Theorem-10.6-conforming inference engine is engineering work and is
deferred to a companion artifact (§13.3).

\textbf{11. Experimental predictions}

We predict the following empirical behavior for an RNR implementation
that follows the constructions of Sections 4 through 6 with a small (a
few megabytes) byte‑level transformer as the underlying predictor.

These predictions are conjectures, not derived results. They are based
on (a) the structural advantages each Type confers in expectation under
its target data class, and (b) coding rates reported for comparable
existing systems on similar corpora. Specific bit‑per‑byte figures
should be read as order‑of‑magnitude estimates rather than precise
targets; the predictions have not been empirically validated and may be
off by 10-30\% in either direction. Empirical validation against NNCP,
PixelCNN++, LZMA at maximum preset, and zstd at maximum level on the
corpora named below is the natural next step.

\textbf{11.1 Type-I, Code12 on text and source code}

Type-I with $E_{12}$-A on English text and source code is predicted
to achieve 1.5 to 2.0 bits per byte before model amortization,
competitive with LZMA at maximum preset (2.1 to 2.5 bits per byte)
and behind NNCP (1.3 to 1.6 bits per byte on similar corpora). The
gap to NNCP is expected to narrow as the Code12 mapping is learned
(variant $E_{12}$-C) rather than fixed.

\textbf{11.2 Type-II on structured positional data}

Type-II on database pages, structured binary records, and instruction
streams is predicted to achieve coding rates 10 to 20\% below LZMA,
because LZMA exploits string repetition that is weak in these classes
while Type-II exploits positional regularity. Against NNCP, the
comparison depends on how well the byte-level autoregressive model
captures positional structure; for data whose joint structure exceeds
what a sequential model can express in moderate context, Type-II is
predicted to dominate.

\textbf{11.3 Type-III-A composition on heterogeneous archives}

Type-III-A scheduling across Type-I, Type-II, and raw fallback is
predicted to dominate any single mode on archives with heterogeneous
content (mixed text/binary, containers, database dumps with
header/payload structure). The per-mode prediction error
$\varepsilon_t$ determines the closeness to the oracle: for
well-calibrated $U$, the composition is within a few percent of the
per-tile oracle best mode.

\textbf{11.4 Type-III-B on ASCII text and instruction streams}

Type-III-B is predicted to recover 0.3 to 0.7 bits per byte over
independent nibble coding on ASCII text, in line with the
$I(H ; L \mid C)$ estimates of Section 6.4. On structured numeric
arrays (e.g., float32 columns), Type-III-B may recover larger gains
if the byte-level dependency is strong across the mantissa structure.

\textbf{11.5 Type-III-C on highly repetitive structured data}

Type-III-C with a $32 \times 32$ token grid on JSON, XML, source
code, and protocol streams is predicted to achieve coding rates
competitive with the best dictionary-based methods (zstd at maximum
level, LZMA at maximum preset), with the advantage growing as the
corpus size grows and the dictionary amortizes.

\textbf{11.6 Negative controls}

On random or encrypted data, all RNR types are predicted to revert to
baseline plus overhead, since no side information reduces residual
uncertainty. On already-compressed data (gzip, xz, zstd output), RNR
is predicted to be slightly worse than passthrough, since the
residual entropy is high and the model cost is wasted.

\textbf{12. Related work}

This work intersects several established research lines.

\textbf{12.1 Source coding with side information}

Slepian and Wolf (1973) established the rate region for separately
encoding correlated sources with joint decoding. Wyner and Ziv (1976)
extended the result to lossy coding with side information at the
decoder. RNR uses the achievability direction of these results with
deterministic side information, an easier setting than the original
distributed coding problem.

\textbf{12.2 Channel coding and syndrome‑based compression}

Pradhan and Ramchandran (2003) introduced DISCUS for compression with
side information using channel codes. Liveris, Xiong, and Georghiades
(2002) used LDPC codes for similar purposes. Type‑I with a structured
E\_12 mapping is in the same spirit: the predictor produces a corrupted
version of a channel code, and the repair stream is essentially syndrome
information.

\textbf{12.3 Enumerative and combinatorial coding}

Cover (1973) gave the enumerative source coding scheme central to
Type‑II. Subsequent works extended enumerative coding to constrained
classes. Theorem 5.3 specializes these to the
partition‑with‑support‑and‑escape setting.

\textbf{12.4 Block transforms}

The Burrows--Wheeler Transform (Burrows and Wheeler 1994) is the
classical example of a reversible reorganization of data that exposes
structure. Type‑II and Type‑III‑B are also reversible reorganizations:
positional partition and byte‑tensor representations.

\textbf{12.5 Neural lossless compression}

Bellard's NNCP (2019) is the closest prior work to Type‑I: a byte‑level
transformer drives an arithmetic coder. PixelRNN (van den Oord et al.,
ICML 2016) and the PixelCNN family --- Conditional PixelCNN (van den
Oord et al., NeurIPS 2016), PixelCNN++ (Salimans et al. 2017), and
PixelSNAIL (Chen et al. 2018) --- are factorized autoregressive image
models that can be viewed as Type‑II coders under a factorized regime;
PixelCNN++ uses discretized logistic mixtures rather than raw
categoricals but the structural reduction is the same. Integer Discrete
Flows (Hoogeboom et al. 2019) and IDF++ (van den Berg et al. 2021)
construct invertible discrete representations whose factorized form is
closely related to Type‑III‑B. L3C (Mentzer et al. 2019) and LC-FDNet
(Rhee et al. 2022) introduce hierarchical priors that prefigure
Type‑III‑A composition.

\textbf{12.6 Bits‑back coding}

Bits‑back coding (Hinton and van Camp 1993; Frey 1997) was made
practical by Townsend, Bird, and Barber (2019) as BB‑ANS, and extended
by Kingma et al. (2019) as Bit‑Swap and by Townsend et al. (2020) as
HiLLoC. These methods losslessly compress data under variational
latent‑variable models and are the appropriate machinery for Type‑III‑C
with continuous latents.

\textbf{12.7 Context mixing}

PAQ (Mahoney) and Cmix (Knoll) represent the upper end of non‑neural
lossless compression via mixture‑of‑experts context mixing. They can be
viewed as hand‑designed Type‑III‑A schedulers in which the entropy field
U is computed from explicit context features rather than predicted by a
neural network.

\textbf{13. Status of open problems}

This section reports the status of every problem the framework opens,
distinguishing problems resolved within the paper from those that are
genuinely open and from those deferred to engineering or empirical work.
The taxonomy follows the structure recommended in earlier framework
reviews: a problem is reported in §13.1 if its answer is a theorem, in
§13.2 if its answer requires additional theoretical or empirical work
that has not been done, and in §13.3 if its answer is a piece of code, a
conformance test, or a policy statement rather than a mathematical
result.

\textbf{13.1 Resolved problems}

The following problems are resolved in this paper. Cross-references
point to the relevant theorems.

\begin{itemize}
\item
  Type-I converse bound --- resolved by Theorem 4.2: the
  achievability bound of Theorem 4.1 matches the conditional
  source-coding lower bound $H(X \mid Y)$ to first order in $N$.
\item
  Computational complexity of optimal $E_{12}$ --- partially
  resolved by Theorem 4.3 (NP-hardness of the varying-parameter
  learned-code labeling problem via Wagner and Corneil's 1990
  tree-to-hypercube subgraph embedding). The fixed Code12 instance
  ($n_0 = 256, m_0 = 12$) is constant-size and hence formally
  polynomial in the input bit-length of the empirical confusion data
  (Remark 4.3a). A sound constant-gap inapproximability proof in
  the Hamming-restricted geometry remains open and is listed in
  §13.2.
\item
  Joint structure on natural data and the Type-II separation regime
  --- partially resolved. Theorem 5.5 is a theorem: the
  multi-information $\mathrm{TC}(D_N)$ is the exact
  information-theoretic quantity governing factorized-versus-joint
  separation. Remark 5.6, which extrapolates this to $\Theta(N)$
  separation under the manifold hypothesis for natural data, is an
  informal heuristic (not a theorem) supported by but not implied by
  the cited empirical estimates of intrinsic dimension $d$; the
  precise scaling of $\mathrm{TC}(D_N)$ for any specific natural
  data class is an empirical question and is also listed in §13.2.
\item
  Tightness of the $2\varepsilon$ mode-selection bound --- resolved
  by Theorem 6.2: the $2\varepsilon$ factor in Theorem 6.1 cannot be
  improved without additional assumptions on the prediction error
  structure.
\item
  Optimal hierarchical decomposition for Type-III-B --- resolved by
  Theorem 6.4: all rooted binary tree decompositions of the 256-byte
  alphabet achieve the same information-theoretic rate
  $H(B \mid C)$; the choice of tree affects only encoder/decoder
  complexity and predictor calibration.
\item
  Bits-back integration into the $L(R \mid Y)$ decomposition ---
  resolved by Theorem 6.6: the three-stream accounting yields
  expected rate equal to the cross-entropy of the data under the
  model marginal plus the variational gap, that is,
  $H_D(X \mid Y) + \mathrm{KL}(D(\cdot \mid Y)\,\Vert\,
  P_X(\cdot \mid Y)) + \mathrm{KL}(q(z \mid X, Y)\,\Vert\,
  P(z \mid X, Y))$; this matches the Slepian-Wolf bound
  $H_D(X \mid Y)$ exactly when both KL terms vanish (model marginal
  matches data, and $q$ matches the true posterior).
\item
  Computational complexity of joint dictionary-and-predictor
  learning --- partially resolved. Theorem 6.7 establishes
  NP-hardness of the \emph{restricted} dictionary-only case via
  reduction from SGP; Theorem 6.8 transfers APX-hardness from the
  SGP literature (Lehman--Shelat 2002; Charikar et al.\ 2005) to the
  same restricted case with the same inherited constant
  $\gamma_{\mathrm{SGP}} > 1$ (we do not pin down a specific
  numerical value). Hardness of the \emph{unrestricted} joint
  problem is conjectured but not proved by the present reductions
  and is listed in \S13.2.
\item
  Bit-exact decoder reproduction for fixed predictors --- resolved
  by Theorem 10.1: a precision-and-determinism specification
  suffices for cross-platform reproducibility, reducing the question
  to a conformance test.
\item
  Bit-exact reproduction under adaptive predictors --- resolved by
  Definition 10.1 (compatible online optimizer) and Theorem 10.2:
  online adaptation via integer-arithmetic optimization is bit-exact
  across implementations, with several published optimizers fitting
  the compatible class.
\item
  Random access (seek-and-read) into a compressed archive with
  output-sensitive cost $O(\log(|X|/K) \cdot \log L(R) + (K + k) \cdot
  \mathrm{cost}(M))$ --- resolved by Definitions 10.2 ($W$-bounded
  predictor) and 10.3 (sync point, including entropy-coder flush)
  and Theorem 10.3, with Theorem 10.4 covering non-$W$-bounded
  predictors via state snapshots. The cost is polylogarithmic in
  archive size $|X|$ only in the parameter regime where the
  sync-point spacing $K$, read window $k$, and $\log L(R)$ are held
  fixed or grow at most polylogarithmically in $|X|$; outside that
  regime (e.g., $K = \Theta(|X|)$ or $k = \Theta(|X|)$) the cost is
  $\Theta(|X| \cdot \mathrm{cost}(M))$. This is the framework's deployment-grade
  distinguishing feature relative to other neural compressors,
  which lack random-access support entirely.
\item
  Block-device and read-only filesystem deployment with
  hardware-target portability --- resolved by Theorem 10.5
  (block-device interface: an OS-level driver serves LBA-indexed
  read requests with the same output-sensitive cost as Theorem 10.3,
  polylogarithmic in $|X|$ only in the same fixed/polylog $K$, $k$,
  $\log L(R)$ parameter regime) and Theorem 10.6 (hardware-target
  portability:
  bit-exact predictor execution on CPU scalar, CPU SIMD, GPU integer
  Tensor Cores, and integer NPUs follows from Theorem 10.1's
  conformance specification, augmented in Theorem 10.6 with explicit
  weight/input byte-serialization conformance). The copy-on-write
  overlay pattern (§10.8) enables read-write semantics using the
  same architecture as existing compressed-filesystem deployments.
\end{itemize}

\textbf{13.2 Remaining mathematical open problems}

A small number of mathematical questions remain genuinely open after
the present paper.

\begin{itemize}
\item
  Optimal candidate-support threshold for semi-non-factorized
  Type-II. The natural marginal-$Q$-prefix heuristic
  ($s$ largest-marginal symbols) is not provably optimal: Remark 5.7
  retracts the corresponding earlier proposition with a concrete
  counterexample at $N = 7, k = 2, s = 2$, demonstrated by the
  verification script in \texttt{scripts/verify/}. A general
  polynomial-time selector or a formal hardness result for support
  optimization is open. The arithmetic-coding benchmark under the
  same $Q$ is a separate, settled result and remains the comparison
  baseline.
\item
  Hardness of the unrestricted joint Type-III-C problem
  $(M, D, R \mid Y)$. Theorems 6.7 and 6.8 establish NP-hardness and
  APX-hardness only for the restricted normal-form
  dictionary-only case under the symbol-count objective
  $L_{\mathrm{sym}}$. The unrestricted problem (model, dictionary,
  repair stream all free, total code length
  $L(M) + L(D) + L(R \mid Y)$) is conjectured to be at least as
  hard, but neither inclusion nor a direct gap-preserving reduction
  is available: the feasible set is strictly larger and the
  objective function differs. A reduction that forces the
  unrestricted optimum to coincide with (or polynomially recover)
  the restricted optimum is open, both for NP-hardness and for
  APX-hardness with any explicit inherited constant.
\item
  Tight asymptotic scaling of $\mathrm{TC}(D_N)$ for specific
  natural-data classes. Theorem 5.5 establishes the
  multi-information $\mathrm{TC}(D_N)$ as the exact separation
  between factorized and joint coding. Remark 5.6's $\Theta(N)$
  extrapolation for low-dimensional natural-data manifolds is an
  informal heuristic; rigorously establishing or refuting
  $\mathrm{TC}(D_N) = \Theta(N)$ for any specific class (e.g.,
  natural images at a given resolution, English text at a given
  vocabulary, database pages under a given schema) is open.
\item
  Construction of efficient enumerative encoders for specific
  natural-data submanifolds. Theorem 5.5 establishes that the
  separation between factorized and joint coding equals
  $\mathrm{TC}(D_N)$, and Remark 5.6 predicts this is $\Theta(N)$
  for low-dimensional natural-data manifolds. Theorem 5.5 itself
  does not identify any particular submanifold or guarantee the
  existence of a polynomial-time membership oracle or rank coder
  for one; for an arbitrary finite distribution $D_N$ a brute-force
  oracle that enumerates the support exists trivially but is not
  efficient. Constructing a polynomial-time membership oracle and
  rank coder for the support of a specific data class --- say,
  database pages under a given schema, or valid x86-64 instruction
  streams --- is open as a constructive question for each class, and
  is not implied by Theorem 5.5.
\item
  Constant-gap inapproximability for $E_{12}$ in the
  Hamming-restricted geometry. The decision version of
  varying-parameter learned-code labeling is NP-hard (Theorem 4.3)
  via an exact YES/NO threshold reduction; no constant-gap reduction
  from any APX-hard source into the Hamming-realizable distance
  geometry is currently known. The Type-III-C side (Theorem 6.8) is
  APX-hard with an inherited constant $\gamma_{\mathrm{SGP}} > 1$
  whose specific numerical value depends on which source-problem
  reduction is used (the literature includes several; we do not pin
  one down here), but the analogous constant for $E_{12}$ is open.
\item
  Practical compatible Adam variants with full convergence
  analysis. Theorem 10.2 covers any optimizer satisfying Definition
  10.1. Integer-arithmetic Adam variants exist (Banner et al. 2018)
  but their convergence behavior relative to floating-point Adam
  under the specific compatibility constraints (deterministic
  integer reciprocal-square-root, fixed update ordering) is
  incompletely characterized. A specific compatible Adam variant
  with a convergence guarantee matching the floating-point
  baseline, on standard problems, is open.
\end{itemize}

\textbf{13.3 Out of scope: engineering deferrals}

The following items are deferred to engineering or empirical work. They
are not theorem-shaped problems and cannot be resolved within the scope
of a theoretical paper. They are listed here only so that a reader can
see what additional work is required for a deployable system.

\begin{itemize}
\item
  Reference implementation conformance. Theorems 10.1 and 10.2 specify
  sufficient conditions for bit-exact reproduction under fixed and
  adaptive predictors. A reference implementation meeting both, plus a
  conformance test suite that verifies cross-platform identity on a
  representative corpus, is required for deployable archives. This is
  software engineering work, not mathematics. Out of scope for this
  paper; the natural artifact is a separate code release accompanied by
  a conformance specification.
\item
  Empirical validation. The quantitative predictions in §11 are
  conjectures based on the structural advantages of each Type and on
  coding rates reported for comparable systems. Validation requires
  running the constructions of §§4--6 against NNCP, PixelCNN++, LZMA,
  zstd, and bits-back baselines on a representative corpus. This is
  experimental work, not mathematics. Out of scope for this paper; the
  natural artifact is a separate experimental paper.
\item
  Intrinsic-dimension estimation for specific data classes. Remark 5.6
  connects Theorem 5.5\textquotesingle s separation to the manifold
  hypothesis. Computing or estimating the intrinsic dimension d of
  specific natural-data classes (database pages of a given schema,
  telemetry under a given protocol, image data of a given modality)
  requires applied data analysis using established estimators
  (Levina-Bickel maximum likelihood, correlation dimension, persistent
  homology). This is empirical work, not mathematics. Out of scope for
  this paper.
\item
  Privacy and information leakage from shipped predictors. A model M
  trained on private data and shipped as part of an archive may leak
  training-set information through memorization or membership inference.
  The trade-off between L(M) reduction (sharing more model) and
  information leakage requires a threat model, an adversary capability
  specification, and an empirical privacy audit. This is a security and
  policy question, not a theorem; out of scope for this paper.
\end{itemize}

\textbf{References}

\textbf{{[}1{]}} Banner, R., Hubara, I., Hoffer, E., and Soudry, D.
(2018). Scalable Methods for 8-bit Training of Neural Networks. NeurIPS.

\textbf{{[}2{]}} Bellard, F. (2019). Lossless Data Compression with
Neural Networks. Technical report and software,
https://bellard.org/nncp/.

\textbf{{[}3{]}} Bernstein, J., Wang, Y.-X., Azizzadenesheli, K., and
Anandkumar, A. (2018). signSGD: Compressed Optimisation for Non-Convex
Problems. ICML.

\textbf{{[}4{]}} Burrows, M., and Wheeler, D. J. (1994). A Block-Sorting
Lossless Data Compression Algorithm. SRC Research Report 124, Digital
Equipment Corporation.

\textbf{{[}5{]}} Carlsson, G. (2009). Topology and Data. Bulletin of the
American Mathematical Society, 46(2), 255-308.

\textbf{{[}6{]}} Carlsson, G., Ishkhanov, T., de Silva, V., and
Zomorodian, A. (2008). On the Local Behavior of Spaces of Natural
Images. International Journal of Computer Vision, 76(1), 1-12.

\textbf{{[}7{]}} Chen, X., Mishra, N., Rohaninejad, M., and Abbeel, P.
(2018). PixelSNAIL: An Improved Autoregressive Generative Model.
International Conference on Machine Learning.

\textbf{{[}8{]}} Charikar, M., Lehman, E., Liu, D., Panigrahy, R.,
Prabhakaran, M., Sahai, A., and Shelat, A. (2005). The Smallest Grammar
Problem. IEEE Transactions on Information Theory, 51(7), 2554-2576.

\textbf{{[}9{]}} Cover, T. M. (1973). Enumerative Source Encoding. IEEE
Transactions on Information Theory, 19(1), 73-77.

\textbf{{[}10{]}} Cover, T. M., and Thomas, J. A. (2006). Elements of
Information Theory, 2nd edition. Wiley.

\textbf{{[}11{]}} Donoho, D. L. (2000). High-Dimensional Data Analysis:
The Curses and Blessings of Dimensionality. AMS Math Challenges Lecture.

\textbf{{[}12{]}} Duda, J. (2009). Asymmetric Numeral Systems.
arXiv:0902.0271.

\textbf{{[}13{]}} Frey, B. J. (1997). Bayesian Networks for Pattern
Classification, Data Compression, and Channel Coding. Ph.D. thesis,
University of Toronto.

\textbf{{[}14{]}} Frey, B. J. (1998). Graphical Models for Machine
Learning and Digital Communication. MIT Press.

\textbf{{[}15{]}} Gallager, R. G. (1962). Low-Density Parity-Check
Codes. IRE Transactions on Information Theory, 8(1), 21-28.

\textbf{{[}16{]}} Hinton, G. E., and van Camp, D. (1993). Keeping Neural
Networks Simple by Minimizing the Description Length of the Weights.
Conference on Learning Theory.

\textbf{{[}17{]}} Hoogeboom, E., Peters, J. W. T., van den Berg, R., and
Welling, M. (2019). Integer Discrete Flows and Lossless Compression.
NeurIPS.

\textbf{{[}18{]}} Huffman, D. A. (1952). A Method for the Construction
of Minimum-Redundancy Codes. Proceedings of the IRE, 40(9), 1098-1101.

\textbf{{[}19{]}} Kingma, F. H., Abbeel, P., and Ho, J. (2019).
Bit-Swap: Recursive Bits-Back Coding for Lossless Compression with
Hierarchical Latent Variables. ICML.

\textbf{{[}20{]}} Knoll, B. (ongoing). Cmix: a lossless data compression
program. https://www.byronknoll.com/cmix.html.

\textbf{{[}21{]}} Levina, E., and Bickel, P. J. (2004). Maximum
Likelihood Estimation of Intrinsic Dimension. NIPS.

\textbf{{[}22{]}} Liveris, A. D., Xiong, Z., and Georghiades, C. N.
(2002). Compression of Binary Sources with Side Information at the
Decoder Using LDPC Codes. IEEE Communications Letters, 6(10), 440-442.

\textbf{{[}23{]}} Mahoney, M. (ongoing). Data Compression Programs and
the PAQ family. http://mattmahoney.net/dc/.

\textbf{{[}24{]}} Mentzer, F., Agustsson, E., Tschannen, M., Timofte,
R., and Van Gool, L. (2019). Practical Full Resolution Learned Lossless
Image Compression. CVPR.

\textbf{{[}25{]}} Pradhan, S. S., and Ramchandran, K. (2003).
Distributed Source Coding Using Syndromes (DISCUS): Design and
Construction. IEEE Transactions on Information Theory, 49(3), 626-643.

\textbf{{[}26{]}} Reed, I. S., and Solomon, G. (1960). Polynomial Codes
over Certain Finite Fields. Journal of SIAM, 8(2), 300-304.

\textbf{{[}27{]}} Rissanen, J. (1978). Modeling by Shortest Data
Description. Automatica, 14(5), 465-471.

\textbf{{[}28{]}} Rissanen, J., and Langdon, G. G. (1979). Arithmetic
Coding. IBM Journal of Research and Development, 23(2), 149-162.

\textbf{{[}29{]}} Salimans, T., Karpathy, A., Chen, X., and Kingma, D.
P. (2017). PixelCNN++: Improving the PixelCNN with Discretized Logistic
Mixture Likelihood and Other Modifications. ICLR.

\textbf{{[}30{]}} Sahni, S., and Gonzalez, T. (1976). P-Complete
Approximation Problems. Journal of the ACM, 23(3), 555-565.

\textbf{{[}31{]}} Shannon, C. E. (1948). A Mathematical Theory of
Communication. Bell System Technical Journal, 27, 379-423 and 623-656.

\textbf{{[}32{]}} Slepian, D., and Wolf, J. K. (1973). Noiseless Coding
of Correlated Information Sources. IEEE Transactions on Information
Theory, 19(4), 471-480.

\textbf{{[}33{]}} Tenenbaum, J. B., de Silva, V., and Langford, J. C.
(2000). A Global Geometric Framework for Nonlinear Dimensionality
Reduction. Science, 290(5500), 2319-2323.

\textbf{{[}34{]}} Townsend, J., Bird, T., and Barber, D. (2019).
Practical Lossless Compression with Latent Variables Using Bits Back
Coding. ICLR.

\textbf{{[}35{]}} Townsend, J., Bird, T., Kunze, J., and Barber, D.
(2020). HiLLoC: Lossless Image Compression with Hierarchical Latent
Variable Models. ICLR.

\textbf{{[}36{]}} van den Berg, R., Gritsenko, A. A., Dehghani, M.,
Sonderby, C. K., and Salimans, T. (2021). IDF++: Analyzing and Improving
Integer Discrete Flows for Lossless Compression. ICLR.

\textbf{{[}37{]}} van den Oord, A., Kalchbrenner, N., and Kavukcuoglu,
K. (2016). Pixel Recurrent Neural Networks. ICML.

\textbf{{[}38{]}} van den Oord, A., Kalchbrenner, N., Vinyals, O.,
Espeholt, L., Graves, A., and Kavukcuoglu, K. (2016). Conditional Image
Generation with PixelCNN Decoders. NeurIPS.

\textbf{{[}39{]}} Rhee, H., Jang, Y. I., Kim, S., and Cho, N. I. (2022).
LC-FDNet: Learned Lossless Image Compression with Frequency
Decomposition Network. CVPR, 6023-6032.

\textbf{{[}40{]}} Witten, I. H., Neal, R. M., and Cleary, J. G. (1987).
Arithmetic Coding for Data Compression. Communications of the ACM,
30(6), 520-540.

\textbf{{[}41{]}} Wu, S., Yan, J., Wang, J., Maxson, A., Xu, K., and He,
K. (2018). Training and Inference with Integers in Deep Neural Networks.
ICLR.

\textbf{{[}42{]}} Wyner, A. D., and Ziv, J. (1976). The Rate-Distortion
Function for Source Coding with Side Information at the Decoder. IEEE
Transactions on Information Theory, 22(1), 1-10.

\textbf{{[}43{]}} Ziv, J., and Lempel, A. (1977). A Universal Algorithm
for Sequential Data Compression. IEEE Transactions on Information
Theory, 23(3), 337-343.

\textbf{{[}44{]}} Ziv, J., and Lempel, A. (1978). Compression of
Individual Sequences via Variable-Rate Coding. IEEE Transactions on
Information Theory, 24(5), 530-536.

\textbf{{[}45{]}} Watanabe, S. (1960). Information Theoretical
Analysis of Multivariate Correlation. IBM Journal of Research and
Development, 4(1), 66-82.

\textbf{{[}46{]}} Studen\'y, M., and Vejnarov\'a, J. (1998). The
Multiinformation Function as a Tool for Measuring Stochastic
Dependence. In: \emph{Learning in Graphical Models}, Jordan, M. I.
(ed.), MIT Press, pp. 261-297.

\textbf{{[}47{]}} Lund, C., and Yannakakis, M. (1994). On the
Hardness of Approximating Minimization Problems. Journal of the ACM,
41(5), 960-981.

\textbf{{[}48{]}} Lehman, E., and Shelat, A. (2002). Approximation
Algorithms for Grammar-Based Compression. Proceedings of the
Thirteenth Annual ACM-SIAM Symposium on Discrete Algorithms (SODA),
205-212.

\end{document}
